\documentclass[11pt,a4paper]{article}
\usepackage[T1]{fontenc}
\usepackage[utf8]{inputenc}
\usepackage{lmodern}
\usepackage[margin=21mm,headheight=14pt]{geometry}
\usepackage{amsmath,amssymb,amsthm,mathtools}
\usepackage{booktabs,longtable,array,enumitem}
\usepackage{xcolor,microtype,fancyhdr,listings}
\usepackage{titlesec}
\usepackage[colorlinks=true,linkcolor=blue!40!black,citecolor=blue!40!black,urlcolor=blue!45!black]{hyperref}
\definecolor{inkblue}{RGB}{24,49,76}
\definecolor{muted}{RGB}{86,99,111}
\definecolor{codebg}{RGB}{247,248,250}
\titleformat{\section}{\Large\bfseries\color{inkblue}}{\thesection}{0.7em}{}
\titleformat{\subsection}{\large\bfseries\color{inkblue}}{\thesubsection}{0.7em}{}
\newtheorem{theorem}{Theorem}[section]
\newtheorem{lemma}[theorem]{Lemma}
\newtheorem{proposition}[theorem]{Proposition}
\newtheorem{corollary}[theorem]{Corollary}
\theoremstyle{definition}
\newtheorem{definition}[theorem]{Definition}
\theoremstyle{remark}
\newtheorem{remark}[theorem]{Remark}
\numberwithin{equation}{section}
\setlist{topsep=5pt,itemsep=3pt,parsep=0pt}
\setlength{\parindent}{0pt}
\setlength{\parskip}{5pt plus 1pt}
\setlength{\emergencystretch}{2em}
\allowdisplaybreaks[2]
\pagestyle{fancy}
\fancyhf{}
\fancyhead[L]{\small\color{muted}Equivalence bridges from the VMRC}
\fancyhead[R]{\small\color{muted}8 November 2026}
\fancyfoot[C]{\small\thepage}
\renewcommand{\headrulewidth}{0.3pt}
\lstset{language=Python,basicstyle=\ttfamily\scriptsize,columns=fullflexible,
  keepspaces=true,breaklines=true,showstringspaces=false,
  backgroundcolor=\color{codebg},frame=single,rulecolor=\color{gray!25},
  xleftmargin=5pt,xrightmargin=5pt,aboveskip=8pt,belowskip=8pt,
  keywordstyle=\color{blue!40!black},commentstyle=\color{green!25!black},
  stringstyle=\color{red!35!black},upquote=true}
\hypersetup{pdftitle={Equivalence Bridges from the 130-Manuscript Corpus},
  pdfauthor={Prepared for Keio},
  pdfsubject={Proofs, exact equivalences, arithmetic obstructions and a positivity boundary},
  pdfkeywords={Bessel, algebraic independence, reciprocal sums, weighted binary expansions, Le Roy, complete monotonicity}}

\begin{document}
\begin{titlepage}
\vspace*{13mm}
{\small\sffamily\color{muted}RESEARCH SYNTHESIS\hfill 8 NOVEMBER 2026\par}
\vspace{13mm}
{\Huge\bfseries\color{inkblue}Equivalence Bridges\par}
\vspace{3mm}
{\LARGE\bfseries\color{inkblue}from the 130-Manuscript Corpus\par}
\vspace{8mm}
{\Large Explicit independence, arithmetic obstructions,\par
and a positivity boundary\par}
\vspace{9mm}
{\large Prepared for Keio\par}
\vspace{8mm}
\hrule
\vspace{7mm}
\textbf{Abstract.}
This note develops seven connections from selected arguments in the
Verifiable Mathematical Reasoning Corpus. Its 28 numbered theorems and
corollaries include a complete algebraic-relation criterion for
logarithm-cancelled Bessel values; density-one independence under rational
height; exact prime-power multiplicity and a unique height-eight
minimizer; a quadratic weighted-binary denominator obstruction;
a number-field extension of the translated reciprocal-tail method;
geometric characterizations of universal integrability; and an exact
contact description of a Le Roy complete-monotonicity boundary.

Three representative conclusions are
\[
 \operatorname{trdeg}_{\mathcal F}\mathcal F(C\kappa)
       =\operatorname{rank}_{\mathbb Q}C,
 \qquad
 \lambda^d\le d[K:\mathbb Q]+1,
\]
and the certified strict lower bound
\[
 \beta_c(1/3,2)>0.428292245841.
\]
Each statement is given with its hypotheses and proof. Exact finite
checks and interval certificates accompany the analytic arguments.
The count includes supporting corollaries and useful classical
equivalences; it is not a claim to 28 independent discoveries or to
28 newly solved conjectures. No historical priority or external
peer-review certification is asserted.
\vfill
{\small\color{muted}
Corpus snapshot: \texttt{c30d40119134a59eda6fa5cc6331f27c8e67d034}.\par
All 130 manuscript files matched the snapshot's Git blob hashes.\par
Mathematical review was targeted at the inputs used in this note.}
\end{titlepage}
\setcounter{page}{2}

\tableofcontents
\clearpage

\section{Sources, results, and proof boundaries}
\label{sec:overview}

\subsection{Verified source scope}

The source repository is the
\href{https://github.com/QEEDDD/Verifiable-Mathematical-Reasoning-Corpus}{Verifiable Mathematical Reasoning Corpus}.
Its \href{https://github.com/QEEDDD/Verifiable-Mathematical-Reasoning-Corpus/tree/c30d40119134a59eda6fa5cc6331f27c8e67d034}{immutable snapshot}
contains 130 numbered \texttt{.tex} manuscripts. On 8 November 2026,
all 130 local file contents were compared with the current repository
tree using Git's blob hash rule and matched. This checks source
identity and coverage of the inventory; it does not verify every
argument in all 130 manuscripts.

The principal inputs inspected for this synthesis are the following.
The bracketed numbers throughout the note refer to corpus manuscript
numbers, while ordinary bibliography citations refer to external works.

\begin{center}
\small
\begin{tabular}{@{}>{\raggedright\arraybackslash}p{0.09\linewidth}p{0.41\linewidth}p{0.43\linewidth}@{}}
\toprule
Source & Subject & Role here\\
\midrule
{[044]} & Euler--Gompertz and Bessel values & Joint E-function independence and affine logarithmic rank\\
{[056]} & Square moduli and vanishing Erd\H{o}s sums & Congruence-constrained periodic coefficients\\
{[080]} & Exponent window for Erd\H{o}s 265 & Positive multilinear tails\\
{[086]} & Three nearby primes in Erd\H{o}s 700 & Lucas digit restrictions\\
{[088]} & Integer-tail analytic attack & Perturbation classes and the P\'olya--Carlson bridge\\
{[112]} & Period obstructions for weighted binary expansions & Multiplicative orders and valuations\\
{[113], [114]} & Polynomial inverses and integration stability & Polynomial primitive towers and geometric obstructions\\
{[116]} & Le Roy positivity beyond infinite divisibility & Inverse Mellin density and the strict positivity gap\\
\bottomrule
\end{tabular}
\end{center}

For Bessel values, Section~\ref{arb:input-audit} reconstructs the
positive-scale joint-independence argument and states separately the
stronger six-coordinate input from [044]. The linear-algebra and
height deductions are then proved here. The discrete, norm, geometric,
and positivity arguments below include the steps on which their new
conclusions depend. Classical specialization, de Rham, moment, and
natural-boundary theorems remain explicitly identified mathematical
inputs, rather than being counted as new proofs.

\subsection{Seven connections and their exact outputs}

\begin{center}
\small
\begin{tabular}{@{}p{0.24\linewidth}p{0.07\linewidth}p{0.62\linewidth}@{}}
\toprule
Theme & Count & Main output\\
\midrule
Bessel independence & 5 & All polynomial relations among cancelled coordinates are exactly the row relations of a rational matrix; explicit infinite independent family and height bounds.\\
Periodic harmonic sums & 5 & Exact attainable zero multiplicities, congruence iff tests, and a unique height-optimal period-64 example.\\
Quadratic binary coding & 5 & Prime-cube denominator obstruction, finite period bounds, and recurrence/natural-boundary equivalences.\\
Binomial gcd values & 1 & Exact four-value distribution, including all multiplicities, in the nearby-prime family.\\
Translated reciprocal sums & 2 & Common number-field degree lower bounds, uniformly across the perturbation class of Erd\H{o}s 263.\\
Algebraic integration & 3 & Obstruction dimension $2g+s-1$ and universal algebraic/elementary integrability criteria.\\
Le Roy positivity & 7 & Largest-contact equivalence, strict finite Hankel hierarchy, certified numerical lower bound, and convex-order comparison.\\
\bottomrule
\end{tabular}
\end{center}

\textbf{Reading the count.}
There are 28 theorem/corollary statements, in addition to auxiliary
lemmas. For example, the six-way equivalence for binary sequences is
a classical analytic consequence used to transfer the new quadratic
denominator obstruction to recurrence theory. The universal integration
criteria likewise make classical geometric machinery explicit.
``Additional'' means an additional conclusion relative to the selected
corpus passages, not a verified worldwide novelty claim.

\textbf{Connections to open questions.}
The number-field result covers the entire relative-perturbation class
asked about in Erd\H{o}s 263, while leaving its singled-out rationality
question unresolved. The Le Roy analysis characterizes the unknown
true boundary without giving its closed form. The Bessel results
give explicit independent combinations while preserving the distinction
between combinations and individual raw second-kind values. The final
section identifies the precise remaining steps, rather than treating
an equivalence itself as a solution of its open endpoint.

\textbf{Literature controls.}
The CM--moment connection is already in Simon \cite{Simon2021}.
E-function specialization is classical \cite{Beukers2006}.
The unperturbed Mahler/Fermat reciprocal values have classical
transcendence results \cite{Coons2011,AdamczewskiFaverjon2023};
they are not the new perturbative conclusion here. Recent cumulant
lifting results \cite{Zhu2026} also make it inappropriate to count a
generic Bell/Hankel reformulation as a new conjecture resolution.


% Arithmetic bridge fragment.  Self-contained except for the common theorem
% environments and the standard packages listed by the parent agent.
\providecommand{\arbQbar}{\overline{\mathbb Q}}
\providecommand{\arbKexp}{K_{\mathrm{exp}}}
\providecommand{\arbkap}{\kappa}

\section{An exact algebraic-relation criterion for Bessel circuits}
\label{arb:circuits}

Put
\[
 k=\arbQbar,\qquad
 \arbKexp=k(e^a:a\in k),\qquad
 \arbkap(t)=\frac{K_0(t)}{I_0(t)}\quad(t>0),
\]
and let
\[
 \mathcal F=\arbKexp\bigl(I_0(t),I_1(t):
                         t\in k\cap\mathbb R_{>0}\bigr).
\]
All assertions over these countably generated fields have their usual
finite-subset meaning.  The harmonic Bessel coordinate is
\[
 \beta(t)=
 \frac{\displaystyle\sum_{n\geq0}
           H_n(t/2)^{2n}/(n!)^2}{I_0(t)},
 \qquad H_0=0.
\]
The connection formula is
\begin{equation}\label{arb:connection}
 \beta(t)=\gamma+\log(t/2)+\arbkap(t).
\end{equation}
The positive-argument part of corpus [044, Theorem
\texttt{cf:joint}] says that the entire family
\(\{\beta(t):t\in k\cap\mathbb R_{>0}\}\) is algebraically
independent over \(\mathcal F\).  We use this precise statement,
not an inference from separate one-point transcendence results.
A verification of the required functional argument is given below.

\begin{theorem}[Complete relation ideal for cancelled coordinates]
\label{arb:relation-ideal}
Let \(t_1,\ldots,t_m\) be distinct positive algebraic numbers.
Let \(C=(c_{ij})\in\mathbb Q^{s\times m}\) satisfy, for every row,
\[
 \sum_{j=1}^m c_{ij}=0,
 \qquad \prod_{j=1}^m t_j^{c_{ij}}=1.
\]
Define
\[
 L_i=\sum_{j=1}^m c_{ij}\arbkap(t_j),\qquad 1\leq i\leq s.
\]
Then
\begin{equation}\label{arb:rank-equivalence}
 \operatorname{trdeg}_{\mathcal F}
       \mathcal F(L_1,\ldots,L_s)=\operatorname{rank}_{\mathbb Q}C.
\end{equation}
More precisely, the kernel of evaluation
\(\mathcal F[X_1,\ldots,X_s]\to\mathbb C\),
\(X_i\mapsto L_i\), is the ideal generated by
\[
 \sum_i u_iX_i,\qquad u\in\mathbb Q^s,\quad u^{\mathsf T}C=0.
\]
In particular, the \(L_i\) are algebraically independent over
\(\mathcal F\) if and only if the rows of \(C\) are linearly
independent over \(\mathbb Q\).
\end{theorem}

\begin{proof}
For positive arguments the multiplicative equation is equivalent to
\(\sum_jc_{ij}\log t_j=0\).  Equation
\eqref{arb:connection} therefore gives
\(L_i=\sum_jc_{ij}\beta(t_j)\).
The \(\beta(t_j)\) are algebraically independent over
\(\mathcal F\), so they can be treated as indeterminates.
Let \(r=\operatorname{rank}C\).  Rational row operations and an
invertible rational change of coordinates on the indeterminates turn
the map \(X\mapsto C\beta\) into
\[
 (X_1,\ldots,X_s)\longmapsto
       (Y_1,\ldots,Y_r,0,\ldots,0).
\]
Its kernel is generated by \(X_{r+1},\ldots,X_s\).
Undoing the changes of coordinates proves both assertions.
\end{proof}

This is a full algebraic-relation classification for the stated class
of expressions.  It does not assert that an arbitrary combination
of second-kind Bessel values satisfies the two cancellation conditions.
For rational arguments the conditions are decidable by exact
prime factorization: use the matrix whose columns are
\((1,(v_p(t_j))_p)\).

\begin{corollary}[A canonical explicit independent family]
\label{arb:canonical}
For every positive rational \(r\) other than \(1\) or a prime, put
\begin{equation}\label{arb:Dr}
 D_r=\arbkap(r)-\arbkap(1)
       -\sum_{p}v_p(r)\bigl(\arbkap(p)-\arbkap(1)\bigr).
\end{equation}
The sum is finite.  The entire family
\[
 \{D_r:r\in\mathbb Q_{>0}\setminus(\{1\}\cup\{\text{primes}\})\}
\]
is algebraically independent over \(\mathcal F\).
In particular every member is transcendental.  Any nonconstant
rational function, with coefficients in \(\mathcal F\), of finitely
many distinct members is transcendental over \(\mathcal F\).
\end{corollary}

\begin{proof}
Unique prime factorization gives
\(\log r=\sum_pv_p(r)\log p\), and the coefficients in
\eqref{arb:Dr} sum to zero.  In any finite list of such rows,
the column indexed by its non-anchor argument \(r\) has coefficient
one in that row and zero in the others.  The rows are therefore
independent.  Apply Theorem~\ref{arb:relation-ideal}.
A rational function field is relatively algebraically closed over
its coefficient field, which proves the last assertion.
\end{proof}

For example, the following four explicitly specified numbers are
jointly algebraically independent, even over \(\mathcal F\):
\begin{align*}
 D_4&=\arbkap(4)+\arbkap(1)-2\arbkap(2),\\
 D_6&=\arbkap(6)+\arbkap(1)-\arbkap(2)-\arbkap(3),\\
 D_9&=\arbkap(9)+\arbkap(1)-2\arbkap(3),\\
 D_{1/2}&=\arbkap(1/2)+\arbkap(2)-2\arbkap(1).
\end{align*}
For illustration,
\[
 (D_4,D_6,D_9)\approx
 (0.2336077763,\ 0.2754838101,\ 0.3183102998).
\]
These numerical values play no role in the proof.

\subsection{Verification of the independence input used here}
\label{arb:input-audit}

Here is a short reconstruction of the positive-scale input of [044].
It explains why the argument requires joint independence and why
rationally dependent scales such as \(1,2,4\) cause no difficulty.
For a fixed nonzero algebraic scale \(t\), write
\[
 f_t(z)=I_0(tz),\qquad
 y_t(z)=\sum_{n\geq0}\frac{H_n(tz/2)^{2n}}{(n!)^2}.
\]
Then
\[
 f_t''+z^{-1}f_t'-t^2f_t=0,
 \qquad y_t''+z^{-1}y_t'-t^2y_t=2f_t'/z,
\]
and \(g_t=y_t-f_t\log z\) is a second homogeneous solution,
with Wronskian
\(f_tg_t'-f_t'g_t=-1/z\).

The homogeneous equation has differential Galois group
\(\mathrm{SL}_2\).  One can check the usual Bessel input directly:
its logarithmic local solution gives nontrivial unipotent monodromy,
and it is irreducible over \(\mathbb C(z)\).
Indeed a rational logarithmic derivative \(r\) would satisfy
\(r'+r^2+r/z-t^2=0\).
At every nonzero finite pole its residue would be one, while it has
no pole at zero.  At infinity it would have expansion
\(r=\pm t-1/(2z)+O(z^{-2})\), contradicting the nonnegative integer
sum of its finite residues.  The classification of algebraic
subgroups of \(\mathrm{SL}_2\), together with the unipotent element,
then excludes every proper irreducible subgroup.

For scales distinct modulo sign the joint group is the product
\(\mathrm{SL}_2^m\).  If a subdirect product were smaller, two of
its adjoint modules would become isomorphic after a finite algebraic
base change.  Their exponential parts at infinity are respectively
\(0,2t_i,-2t_i\) and \(0,2t_j,-2t_j\).
An invertible formal intertwiner forces equality of these sets,
and hence \(t_i=\pm t_j\), a contradiction.
This is the argument in [044, Lemma \texttt{cf:product}].

Adjoin any finite collection of exponential functions with algebraic
exponents and \(\log z\).  Their Picard--Vessiot group is solvable,
whereas \(\mathrm{SL}_2^m\) has no nontrivial solvable quotient.
The groups therefore form a direct product.
Since the full Bessel solution field is recovered from
\(f_t,f_t',y_t\) after adjoining \(\log z\), dimension counting
gives functional algebraic independence of these \(3m\) coordinates
over the exponential field.
All \(f_t,f_t',y_t,y_t'\) are E-functions in a first-order system
whose only finite singularity is zero.  Ordinary-point
Siegel--Shidlovskii specialization at \(z=1\) applies
\cite{Beukers2006}.
The identity
\[
 f_t'y_t-f_ty_t'=(1-f_t^2)/z
\]
recovers \(y_t'(1)\) from the other three values, since
\(I_0(t)>0\) for positive \(t\).
Lindemann--Weierstrass gives the expected dimension of the
exponential value field.  Thus the \(3m\) values are jointly
independent over each finite exponential coefficient field.
Passing to \(I_0,I_1,\beta\) is birational.  Every proposed relation
over \(\mathcal F\) uses only finitely many additional scales and
exponentials, so the same argument proves the claimed infinite-base
independence.  The specialization input is classical; the scaled
product argument is explicitly checked here and in [044].

\section{From integer arguments to rational height and number fields}
\label{arb:height}

For a finite set \(T=\{t_1,\ldots,t_m\}\) of distinct positive
algebraic numbers, define
\[
 W_T=\{c\in\mathbb Q^m:
             \sum_jc_j=0,\ \sum_jc_j\log t_j=0\},
 \qquad d_T=m-\dim_{\mathbb Q}W_T.
\]
Theorem~\ref{arb:relation-ideal} gives \(m-d_T\) independent
linear combinations of \(\arbkap(t_j)\) over \(\mathcal F\).
Consequently
\begin{equation}\label{arb:Krank}
 \operatorname{trdeg}_{\arbKexp}
    \arbKexp(K_0(t):t\in T)\geq |T|-d_T.
\end{equation}
This is the affine-rank theorem in [044, Theorem \texttt{kd:rank}],
now with an exact relation-ideal formulation for its cancelled outputs.

The stronger six-coordinate input of [044, Theorem
\texttt{jden:rank}] concerns
\begin{align*}
 Z_\lambda={}&\bigl(E_1^{\rm tail}(\lambda),
 K_0(2\lambda),K_1(2\lambda),
 \pi Y_0(2\lambda),\pi Y_1(2\lambda),\mathcal L(\lambda)\bigr),\\
 E_1^{\rm tail}(\lambda)={}&
     \int_\lambda^\infty e^{-u}\,du/u,\\
 \mathcal L(\lambda)={}&
     \sum_{n\geq1}\frac{(-1)^{n-1}\lambda^n}{n(n!)^2}
          -\gamma-\log\lambda.
\end{align*}
It states
\begin{equation}\label{arb:six-input}
 \operatorname{trdeg}_{\arbKexp}
   \arbKexp(Z_\lambda:\lambda\in T)\geq 6|T|-d_T.
\end{equation}
For clarity, this uses more than the positive-scale Bessel input:
it also uses the ordinary-Bessel scales, the entire exponential
integral, and the primitive nonexactness lemma
\texttt{msprim:nonexact} in [044].  The latter reduces to the impossibility
of a rational solution of
\[
 -b''+b'/z-b/z^2+t^2b=1/z.
\]
Nonzero finite poles are impossible.  At zero a nonpositive lowest
Laurent exponent produces an uncancellable pole of order at least
two; a positive lowest exponent produces no \(z^{-1}\) term.
This verifies the critical rational obstruction in that source.
The algebraic recovery in [044, \texttt{jden:rank}] uses eight
independent E-function coordinates, two additional first-kind
coordinates, and the common logarithmic family
\(\gamma+\log\lambda\), of rank at most \(d_T\).
Thus \(8|T|\leq \operatorname{trdeg}(Z_T)+2|T|+d_T\),
which is \eqref{arb:six-input}.

\begin{theorem}[Rational-height density and simultaneous independence]
\label{arb:rationalheight}
Let
\[
 \mathcal R_N=\{a/b:1\leq a,b\leq N,\ \gcd(a,b)=1\},
 \qquad M_N=|\mathcal R_N|=2\sum_{n=1}^N\varphi(n)-1.
\]
Then
\begin{align}
 \operatorname{trdeg}_{\arbKexp}
    \arbKexp(K_0(r):r\in\mathcal R_N)
       &\geq M_N-\pi(N)-1,\label{arb:rationalK}\\
 \operatorname{trdeg}_{\arbKexp}
    \arbKexp(Z_{r/2}:r\in\mathcal R_N)
       &\geq 6M_N-\pi(N)-1.\label{arb:rational6}
\end{align}
There is a single subset \(\mathcal S\subset\mathbb Q_{>0}\)
with
\[
 |\mathcal S\cap\mathcal R_N|\geq M_N-\pi(N)-1
 \qquad(N\geq1)
\]
such that all six coordinates of \(Z_{r/2}\), as
\(r\) ranges over \(\mathcal S\), are jointly algebraically
independent over \(\arbKexp\).
At most \(\pi(N)+1\) arguments \(r\in\mathcal R_N\) can have
any one of
\[
 E_1^{\rm tail}(r/2),\quad \mathcal L(r/2),\quad
 K_j(r),\quad \pi Y_j(r)\quad(j\geq0)
\]
algebraic over \(\arbKexp\).
\end{theorem}

\begin{proof}
The logarithms of the rationals in \(\mathcal R_N\) span exactly
the rational space generated by \(\log p\) for primes \(p\leq N\).
The augmented exponent vectors of \(1\) and of these primes are
independent.  Hence \(d_{\mathcal R_N}=\pi(N)+1\).
Common multiplication by \(1/2\) preserves affine logarithmic
rank, so \eqref{arb:Krank} and \eqref{arb:six-input} give the two
bounds.  The counting formula for \(M_N\) follows by counting the
coprime pairs with \(a\leq b\), and then reflecting across the
diagonal; only \((1,1)\) is on the diagonal.

Enumerate positive rationals by nondecreasing \(\max(a,b)\), and
process the six coordinates in each block in a fixed order.  Keep a
coordinate exactly when it is transcendental over the field generated
by the preceding kept coordinates and \(\arbKexp\).
At the end of \(\mathcal R_N\), at most \(\pi(N)+1\) coordinates
have been rejected by \eqref{arb:rational6}.  Delete every block
containing a rejected coordinate.  At most the same number of blocks
is deleted, and all remaining coordinates are jointly independent.

For the final statement, let \(T\) be the set of bad arguments.
The standard integer-order recurrences express every
\(K_j(r)\), respectively \(\pi Y_j(r)\), as a nonzero
algebraic linear form in the corresponding orders zero and one.
Thus each bad block has a polynomial relation and transcendence
degree at most five.  If \(m=|T|\), its total rank is at most
\(5m\), whereas \eqref{arb:six-input} gives at least
\(6m-d_T\).  Therefore \(m\leq d_T\leq\pi(N)+1\).
\end{proof}

The independent set is existential; this proof does not provide a
membership test or identify an individual algebraic or transcendental
raw second-kind value.  It does, however, give a much sparser
exception proportion under rational height:
\[
 M_N=\frac{6}{\pi^2}N^2+O(N\log N),\qquad
 \frac{\pi(N)+1}{M_N}=O\!\left(\frac1{N\log N}\right).
\]
For example,
\begin{center}
\begin{tabular}{rrrrr}
\toprule
\(N\)&\(M_N\)&block deletion bound&retained blocks&six-coordinate rank bound\\
\midrule
10&63&5&58&373\\
100&6087&26&6061&36496\\
1000&608383&169&608214&3650129\\
\bottomrule
\end{tabular}
\end{center}
The numerical table uses only exact totient and prime counts.

\begin{corollary}[Stability under finitely many additional constants]
\label{arb:base-extension}
Let \(\mathcal B\supset\arbKexp\) be a subfield of \(\mathbb C\)
with finite \(d=\operatorname{trdeg}_{\arbKexp}\mathcal B\).
Every bound in Theorem~\ref{arb:rationalheight} remains valid over
\(\mathcal B\) after adding \(d\) to its deletion or exception
budget, or subtracting \(d\) from its rank bound.
In particular, adjoining any fixed finite list of unresolved
constants preserves the density-one conclusion.
\end{corollary}

\begin{proof}
For any finite tuple \(x\),
\[
 \operatorname{trdeg}_{\mathcal B}\mathcal B(x)
 \geq\operatorname{trdeg}_{\arbKexp}\arbKexp(x)-d.
\]
Apply this to the rank inequalities, then repeat the same greedy
argument and bad-block count.  A list of \(d_0\) constants gives
\(d\leq d_0\), irrespective of their presently unknown arithmetic
natures.
\end{proof}

\begin{theorem}[Density one in a fixed real-embedded number field]
\label{arb:numberfield}
Let \(K\) be a number field of degree \(d_K\) with a fixed real
embedding \(\sigma_0\), and let
\(u_K=r_1+r_2-1\) be its unit rank.  Identify its elements with
their images under \(\sigma_0\), and define
\[
 \mathcal A_X=\{\alpha\in\mathcal O_K:
       \sigma_0(\alpha)>0,\ |\sigma(\alpha)|\leq X
       \text{ for every embedding }\sigma:K\hookrightarrow\mathbb C\}.
\]
Put \(M_K(X)=|\mathcal A_X|\) and
\[
 B_K(X)=d_K\pi(X^{d_K})+u_K+1.
\]
Then
\begin{align*}
 \operatorname{trdeg}_{\arbKexp}
  \arbKexp(K_0(\alpha):\alpha\in\mathcal A_X)
    &\geq M_K(X)-B_K(X),\\
 \operatorname{trdeg}_{\arbKexp}
  \arbKexp(Z_{\alpha/2}:\alpha\in\mathcal A_X)
    &\geq 6M_K(X)-B_K(X).
\end{align*}
There is a single set of retained arguments having at least
\(M_K(X)-B_K(X)\) members in every \(\mathcal A_X\), on which
all six coordinates are jointly independent.  The retained set has
relative density one in these expanding regions.  The same conclusion
survives a fixed finite-transcendence-degree coefficient extension.
\end{theorem}

\begin{proof}
For \(\alpha\in\mathcal A_X\),
\(0<|N_{K/\mathbb Q}(\alpha)|\leq X^{d_K}\).
Every prime ideal dividing \((\alpha)\) lies over a rational prime
at most \(X^{d_K}\).  There are at most
\(d_K\pi(X^{d_K})\) such prime ideals.  The valuation map from
the multiplicative group supported on these primes to the free
abelian group on the primes has kernel \(\mathcal O_K^\times\).
Consequently its rank is at most
\(d_K\pi(X^{d_K})+u_K\).
No assumption that the prime ideals are principal is needed.
Thus \(d_{\mathcal A_X}\leq B_K(X)\), and the rank and greedy
arguments apply exactly as before.

The Minkowski embedding of \(\mathcal O_K\) is a lattice of
covolume \(2^{-r_2}\sqrt{|\Delta_K|}\).  The region consists of a
half-box in the real coordinates and disks in the complex
coordinates.  Elementary lattice counting gives
\[
 M_K(X)=
 \frac{2^{r_1+r_2-1}\pi^{r_2}}{\sqrt{|\Delta_K|}}X^{d_K}
       +O_K(X^{d_K-1}).
\]
On the other hand \(B_K(X)=O_K(X^{d_K}/\log X)\).
Their ratio tends to zero.  Finally use
Corollary~\ref{arb:base-extension}.
\end{proof}

\section{Prime-power coefficient congruences and high-order harmonic zeros}
\label{arb:primepowers}

The congruence classification of corpus [056] concerns the existence
of rational Okada vectors.  The following statements refine the
prime-power construction by measuring the order of vanishing of the
entire periodic Dirichlet series at \(s=1\).  Their scope is explicit:
the first theorem classifies a separated divisor-layer family, and
the second classifies functions constant on each gcd layer.  Neither
is a classification of binary sign functions at arbitrary square
moduli.

Let \(\ell\) and \(p\) be primes with \(p\equiv1\pmod\ell\).
Define the \(p\)-periodic integer function
\[
 u(0)=0,\qquad u(1)=2-p,\qquad
 u(a)=1\quad(2\leq a\leq p-1).
\]
Its period sum is zero and all its nonzero-residue values are one
modulo \(\ell\).  Furthermore
\begin{equation}\label{arb:unonzero}
 L(1,u)<0.
\end{equation}
Indeed, each length-\(p\) block is the sum of \(p-2\) positive
reciprocals with denominators larger than the denominator of the
negative term of magnitude \(p-2\).  Thus each block is strictly
negative, and their sum converges.

\begin{theorem}[Exact maximal zero order in separated prime-power layers]
\label{arb:separated}
Let \(e\geq1\).  For rational
\(c_0,\ldots,c_{e-1}\in1+\ell\mathbb Z_{(\ell)}\), define the
\(p^e\)-periodic function
\[
 f(n)=
 \begin{cases}
 c_tu(n/p^t),&v_p(n)=t<e,\\
 0,&p^e\mid n.
 \end{cases}
\]
Then its period sum is zero, and
\[
 L(s,f)=C(p^{-s})L(s,u),\qquad
 C(T)=\sum_{t=0}^{e-1}c_tT^t.
\]
The maximum possible order of vanishing at \(s=1\), as the
\(c_t\) vary, is exactly
\begin{equation}\label{arb:maxorder}
             \ell^{v_\ell(e)}-1.
\end{equation}
For every integer \(r\geq1\), the following are equivalent:
\begin{enumerate}
\item At least one such \(f\) satisfies \(L^{(j)}(1,f)=0\) for
      \(0\leq j<r\).
\item \(\ell^{v_\ell(e)}\geq r+1\).
\item \(\ell^{\lceil\log_\ell(r+1)\rceil}\) divides \(e\).
\end{enumerate}
The maximal order is attained with integer coefficients and exact
fundamental period \(p^e\).
\end{theorem}

\begin{proof}
The Dirichlet-series identity follows by partitioning the summation
according to \(v_p(n)\).  Each divisor layer has zero mean because
\(u\) has zero mean.  By \eqref{arb:unonzero}, the zero order is
exactly the root multiplicity of \(C(T)\) at \(T=1/p\).

Modulo \(\ell\), one has
\(\overline C(T)=S_e(T)=1+T+\cdots+T^{e-1}\), and \(1/p\)
reduces to one.  Write \(e=Dm\), where
\(D=\ell^{v_\ell(e)}\) and \(\ell\nmid m\).
In characteristic \(\ell\),
\[
 S_e(T)=\frac{(T^m-1)^D}{T-1}
\]
has a root of exact multiplicity \(D-1\) at one.
Since \(p\) is an \(\ell\)-adic unit, any factor
\((T-1/p)^r\) of \(C\) over the rationals can be divided out
over \(\mathbb Z_{(\ell)}\); reduction therefore gives
\(r\leq D-1\).

Equality is attained by
\begin{equation}\label{arb:attainer}
 C(T)=(1-pT)^{D-1}\sum_{j=0}^{m-1}(pT)^{Dj}.
\end{equation}
Its reduction is
\((1+T+\cdots+T^{D-1})\sum_{j=0}^{m-1}T^{Dj}=S_e(T)\),
so every coefficient is one modulo \(\ell\).
The second factor has value \(m\ne0\) at \(T=1/p\).
This proves the exact order and the equivalences.
No off-zero coefficient of \(f\) vanishes, since it is one
modulo \(\ell\).  As \(p^e\) is a period and its zero set is
exactly the multiples of \(p^e\), its fundamental period is
\(p^e\).
\end{proof}

\begin{corollary}[Three logarithmic harmonic cancellations]
\label{arb:81}
There is an odd-integer-valued periodic function of exact period
\(81\), zero exactly at multiples of \(81\), such that
\[
 \sum_{n\geq1}\frac{f(n)}n
 =\sum_{n\geq1}\frac{f(n)\log n}n
 =\sum_{n\geq1}\frac{f(n)(\log n)^2}n=0.
\]
One explicit choice is the construction above with \(p=3\),
\(e=4\), \(u(1)=-1,u(2)=1\), and
\[
 (c_0,c_1,c_2,c_3)=(1,-9,27,-27).
\]
It satisfies
\[
 L(s,f)=-\bigl(1-3^{1-s}\bigr)^3L(s,\chi_{-3}),
 \qquad
 L'''(1,f)=-\frac{2\pi}{\sqrt3}(\log3)^3\ne0.
\]
More generally, for any prescribed \(r\), a suitable
\(p^{2^k}\) supports an odd-integer-valued function whose first
\(r\) logarithm-weighted harmonic sums vanish.
\end{corollary}

\begin{proof}
Here \(u=-\chi_{-3}\), and \(C(T)=(1-3T)^3\).
Use \(L(1,\chi_{-3})=\pi/(3\sqrt3)\) and differentiate the
factor of exact order three.  For a zero-mean periodic function,
the Dirichlet series and its derivatives converge in the half-plane
\(\Re s>0\) after the usual partial summation; in particular the
displayed logarithm-weighted series converge and equal the indicated
derivatives up to the sign \((-1)^j\).
Choose \(2^k-1\geq r\) for the last assertion.
\end{proof}

\begin{theorem}[Exact gcd-layer criterion for arbitrary congruence modulus]
\label{arb:gcdlayers}
Let \(p\) be prime, \(M\geq2\) with \(\gcd(p,M)=1\), and
let \(R_M\) be the rational numbers with denominator coprime to
\(M\).  Fix \(e\geq1\) and \(r\geq0\).
There are coefficients \(c_0,\ldots,c_{e-1}\in1+MR_M\)
for which the function
\[
 f(n)=c_{v_p(n)}\quad(v_p(n)<e),\qquad f(n)=0\quad(p^e\mid n)
\]
has zero mean and a zero of order at least \(r\) at \(s=1\)
if and only if
\begin{equation}\label{arb:hasse}
 \sum_{t=j}^{e-1}\binom tj p^{-(t-j)}\equiv0\pmod{MR_M}
                  \quad(0\leq j\leq r).
\end{equation}
An empty sum is zero, but if \(r\geq e-1\) the test with
\(j=e-1\) is one and hence prevents existence.
The coefficients can be chosen to be integers with the same
congruences.
\end{theorem}

\begin{proof}
For \(\Re s>1\),
\[
 L(s,f)=(1-p^{-s})\zeta(s)C(p^{-s}),
 \qquad C(T)=\sum_{t=0}^{e-1}c_tT^t.
\]
The pole of \(\zeta\) has order one and the other Euler factor
is nonzero at one.  Thus zero mean and zero order at least \(r\)
are equivalent to divisibility of \(C\) by
\((T-1/p)^{r+1}\).
Because \(C\equiv S_e\pmod M\), this implies the vanishing
of the first \(r+1\) Hasse derivatives of \(S_e\) at \(1/p\),
which is precisely \eqref{arb:hasse}.

Conversely subtract from \(S_e(T)\) its Taylor polynomial of
degree \(r\) at \(1/p\).  Under \eqref{arb:hasse}, every
coefficient of the Taylor polynomial being subtracted lies in \(MR_M\), so the resulting
polynomial \(C\) is congruent to \(S_e\) and is divisible by
\((T-1/p)^{r+1}\).
Its coefficients have denominators prime to \(M\).
Multiply by a common denominator chosen to be one modulo \(M\)
to obtain integer coefficients, preserving the congruences and
the zero order.
\end{proof}

\begin{corollary}[Closed form for prime congruence moduli]
\label{arb:gcdprime}
Take \(M=\ell\ne p\) prime and put
\(a=\operatorname{ord}_\ell(p)\) and
\(D=\ell^{v_\ell(e)}\).
The largest root multiplicity at \(T=1/p\) attainable by a
polynomial \(C\equiv S_e\pmod\ell\) is
\[
 \mu=
 \begin{cases}
 0,&a\nmid e,\\
 D,&a\mid e,\ a>1,\\
 D-1,&a=1.
 \end{cases}
\]
Accordingly, a convergent gcd-layer series exists exactly when
\(\mu\geq1\), and its maximal zero order at one is
\(\mu-1\).  In particular, for odd \(p\) and odd coefficients,
the maximal order in this gcd-layer class is
\(2^{v_2(e)}-2\), when a convergent series exists.
\end{corollary}

\begin{proof}
In characteristic \(\ell\), write \(e=Dm\) with
\(\ell\nmid m\), and factor
\(T^e-1=(T^m-1)^D\).
The point \(p^{-1}\) is a root precisely when \(a\mid e\).
Its multiplicity in \(T^e-1\) is then \(D\), with one removed
by division by \(T-1\) exactly when \(a=1\).
For attainability, divide \(S_e\) in \(\mathbb F_\ell[T]\)
by the maximal power of \(1-\overline pT\), lift the quotient
to an integer polynomial, and multiply by \((1-pT)^\mu\).
The lifted quotient is nonzero at \(1/p\), as its reduction is
nonzero there.  Hence the multiplicity is exactly \(\mu\).
\end{proof}

\begin{corollary}[A height-optimal double zero at period 64]
\label{arb:64-optimal}
Define \(f(n)=c_{v_2(n)}\) when \(v_2(n)<6\), and zero when
\(64\mid n\), using
\[
 (c_0,c_1,c_2,c_3,c_4,c_5)=(1,-5,7,-2,4,-8).
\]
All its off-zero values are one modulo three, its fundamental period
is 64, and
\[
 L(s,f)=(1-2^{-3s})\zeta(s)(1-2^{1-s})^3.
\]
Consequently
\[
 L(1,f)=L'(1,f)=0,\qquad
 L''(1,f)=\frac74(\log2)^3\ne0.
\]
Its height \(\max_n|f(n)|=8\) is the least possible among integer
gcd-layer functions of period 64, congruent to one modulo three
off zero, whose Dirichlet series vanishes to order at least two
at one. The displayed coefficient vector is the unique minimizer in
that class.
\end{corollary}

\begin{proof}
The coefficient polynomial factors as
\[
 1-5T+7T^2-2T^3+4T^4-8T^5
     =(1-2T)^3(1+T+T^2).
\]
Use the identity in Theorem~\ref{arb:gcdlayers} and
\((1-T)(1+T+T^2)=1-T^3\).
The factor \(1-2^{1-s}\) has a simple zero with leading
coefficient \(\log2\), while \(\zeta(s)\) has residue one.
The coefficient of \((s-1)^2\) is therefore
\((7/8)(\log2)^3\), giving the stated second derivative.

For optimality, any competing integer polynomial of degree five
must have the factor \((1-2T)^3\).
That factor is primitive in \(\mathbb Z[T]\), so Gauss's lemma
makes its quotient an integer polynomial.  Its leading coefficient
\(c_5\) is a nonzero multiple of eight; it cannot be zero since
\(c_5\equiv1\pmod3\).  The height is thus at least eight,
and the displayed example attains it.

For uniqueness at height eight, write
\(C(T)=(1-2T)^3(A+BT+ET^2)\), where \(A,B,E\in\mathbb Z\).
The conditions on \(c_5=-8E\) force \(E=1\).
Then \(c_4=12-8B\), its height bound and its congruence force \(B=1\).
Finally \(c_0=A\equiv1\pmod3\) and \(|c_3|=|6-8A|\le8\)
force \(A=1\). This gives exactly the stated vector.
\end{proof}

These two prime-power theorems expose a genuine distinction.
The nonconstant zero-mean layer seed in
Theorem~\ref{arb:separated} removes the zeta pole before the
coefficient polynomial is used.  Constant gcd layers must spend
one additional root on that pole.  The factor-of-two discrepancy
between an arbitrary odd Okada vector and a gcd-invariant one at
small prime-power exponents is therefore an exact algebraic effect.
The constructions have odd or prescribed-congruence coefficients,
not coefficients restricted to \(\{\pm1\}\); they do not solve
the remaining square-modulus sign problem.


\section{Quadratic weighted binary expansions}
\label{disc:quadratic}

The valuation argument of corpus [112] concerns the weight $n$ in a
binary expansion.  The following extension concerns the weight $n^2$.
All valuation and representation arguments needed here are included.
The analytic consequence uses the classical P\'olya--Carlson theorem,
also used in corpus [088].

For binary digits $\varepsilon_n\in\{0,1\}$, put
\[
 X_2(\varepsilon)=\sum_{n\ge1}\frac{n^2\varepsilon_n}{2^n}.
\]
For an odd prime $p$, write
\[
 t_p=\operatorname{ord}_p(2),\qquad
 w_p=v_p(2^{t_p}-1)\ge1.
\]

\begin{theorem}[Exact quadratic denominator amplification]
\label{disc:amplification}
Suppose the digits of $X_2$ have a nonconstant periodic tail, beginning
after position $N$, of least period $r$.  Let $d>1$ be the reduced
denominator of the ordinary binary value of this periodic tail, and let
$Q$ be the odd part of the reduced denominator of $X_2$.  If $p\mid d$
and
\[
 a=v_p(d)>b=v_p(r),
\]
then
\[
 \boxed{v_p(Q)=a+2w_p.}
\]
In particular, $Q$ is divisible by the cube of some odd prime.
\end{theorem}

\begin{proof}
Write the periodic word as $e_1,\ldots,e_r$, and set
\[
 M=2^r-1,\quad
 S=\sum_{i=1}^r e_i2^{r-i},\quad
 T=\sum_{i=1}^r i e_i2^{r-i},\quad
 U=\sum_{i=1}^r i^2 e_i2^{r-i}.
\]
Summing over the positions $i+kr$ gives the three tail values
\begin{align*}
 y&=\sum_{i\ge1}\frac{e_i}{2^i}=\frac{S}{M},\\
 z&=\sum_{i\ge1}\frac{i e_i}{2^i}=\frac{T}{M}+\frac{rS}{M^2},\\
 w&=\sum_{i\ge1}\frac{i^2e_i}{2^i}
   =\frac{U}{M}+\frac{2rT}{M^2}
      +\frac{r^2S(M+2)}{M^3}.
\end{align*}
Here the digits on the left are extended periodically.  The last
formula follows from
$\sum_{k\ge0}k^2q^k=q(1+q)/(1-q)^3$ with $q=2^{-r}$.
The prefix contributes an integer $I$ after multiplication by $2^N$,
so
\begin{equation}\label{disc:prefix}
 2^NX_2=I+N^2y+2Nz+w.
\end{equation}

The ordinary radix-expansion theorem gives $r=\operatorname{ord}_d(2)$.
Put $k=v_p(M)$.  The lifting formula at an odd prime gives
$k=w_p+b$, while reduction of $S/M$ gives $v_p(S)=k-a$.
Since $M+2$ is a $p$-adic unit, the last term in $w$ has exact valuation
\[
 2b+(k-a)-3k=-a-2w_p.
\]
The middle term has valuation at least
$b-2k=-b-2w_p$, and the first at least $-k=-b-w_p$.
Both bounds are strictly larger because $a>b$.  Hence
$v_p(w)=-a-2w_p$.  Similarly, the displayed formula for $z$ gives
$v_p(z)=-a-w_p$, while $v_p(y)=-a$.  Thus no other term in
\eqref{disc:prefix} can cancel the term of least valuation.  This proves
the exact formula for $Q$.

To obtain the final assertion, choose the largest prime $p\mid d$.
For an odd prime $q$ and an integer $c\ge1$, the order formula is
\[
 \operatorname{ord}_{q^c}(2)
 =t_q q^{\max(0,c-w_q)}.
\]
Taking the least common multiple of these orders over $q\mid d$ shows
that $b=\max(0,a-w_p)<a$: a smaller prime $q<p$ contributes no factor
$p$, because $t_q\mid q-1<p$.  The exact formula therefore gives
$v_p(Q)=a+2w_p\ge3$.
\end{proof}

\begin{corollary}[Aperiodicity for cube-free odd denominators]
\label{disc:cubefree}
Every $x\in[0,6]$ has a representation
\[
 x=\sum_{n\ge1}\frac{n^2\varepsilon_n}{2^n},
 \qquad \varepsilon_n\in\{0,1\}.
\]
If $x$ is rational and the odd part $Q>1$ of its reduced denominator
is cube-free, then every such representation is aperiodic after every
finite prefix.  This applies, for example, to both $1/3$ and $1/9$.
The prime-cube obstruction is sharp: the periodic word $10$ gives
\[
 \sum_{k\ge0}\frac{(2k+1)^2}{2^{2k+1}}=\frac{82}{27}.
\]
\end{corollary}

\begin{proof}
Let $u_n=n^2/2^n$.  Its future tail is
\[
 T_n=\sum_{j>n}u_j=\frac{n^2+4n+6}{2^n},
\]
so $u_n\le T_n$ for every $n$ and the total sum is $6$.
Given a remainder in $[0,u_n+T_n]$, choose the next digit so that the
new remainder lies in $[0,T_n]$.  This is possible because
$[0,T_n]\cup[u_n,u_n+T_n]=[0,u_n+T_n]$.  Since $T_n\to0$, the
construction proves existence for every $x\in[0,6]$.

A nonconstant periodic tail contradicts
Theorem~\ref{disc:amplification}.  An eventually zero tail gives a
dyadic rational, and an eventually one tail does too, by the formula
for $T_n$.  Neither is possible when $Q>1$.  Finally, the word $10$
has $r=2$, $M=3$, and $S=T=U=2$, so its quadratic weighted value is
$2/3+8/9+40/27=82/27$.
\end{proof}

\begin{theorem}[Finite bound on all eventual periods]
\label{disc:periodbound}
For an odd positive integer $Q$, let $\mathcal P_Q^{(2)}$ be the
smallest set of odd primes containing every $p$ with
$v_p(Q)>2w_p$, and containing every odd prime divisor of $t_q$
whenever it contains $q$.  This set is finite.  Define
\[
 e_p=\max\left\{0,\ v_p(Q)-2w_p,\
       \max_{\substack{q\in\mathcal P_Q^{(2)}\\q>p}}v_p(t_q)\right\},
 \qquad D_Q^{(2)}=\prod_{p\in\mathcal P_Q^{(2)}}p^{e_p},
\]
with empty maxima and products interpreted as $0$ and $1$.
For every nonconstant eventual tail of a quadratic weighted binary
representation with odd denominator $Q$, its ordinary binary
denominator $d$ and least period $r$ satisfy
\[
 d\mid D_Q^{(2)},\qquad
 r\mid\operatorname{ord}_{D_Q^{(2)}}(2).
\]
In particular, such a tail requires $D_Q^{(2)}>1$.
\end{theorem}

\begin{proof}
Every newly added prime is smaller than its parent because
$t_q\mid q-1$, proving finiteness.  Write $a_p=v_p(d)$ and
$b_p=v_p(r)$.  If $a_p>b_p$, Theorem~\ref{disc:amplification}
gives $a_p=v_p(Q)-2w_p$.  Otherwise $a_p\le b_p$, and the order
formula in the preceding proof gives
\[
 b_p=\max\left\{\max(0,a_p-w_p),\
          \max_{\substack{q\mid d\\q>p}}v_p(t_q)\right\}.
\]
Its first entry is strictly smaller than $a_p$.  Thus a larger
$q\mid d$ satisfies $a_p\le v_p(t_q)$.  Process the prime factors
of $d$ in decreasing order.  Each is either an initial seed with the
stated exponent bound, or is introduced by an already processed
larger prime and again has exponent at most $e_p$.  This proves
$d\mid D_Q^{(2)}$.  Divisibility of multiplicative orders gives
the assertion about $r$.
\end{proof}

This is a necessary period bound, not a sufficient condition for the
existence of a periodic representation of a specified numerator.
For $Q=27$, it gives $D_Q^{(2)}=3$ and $r\mid2$.
For a fully finite example, trial division through $31$ proves that
$1093$ is prime, and exact modular exponentiation gives
\[
\begin{array}{c|rrrr}
m&364&182&52&28\\\hline
2^m\bmod1093&1&1092&27&121.
\end{array}
\]
Since the prime divisors of $364$ are $2,7,13$, these congruences
prove $\operatorname{ord}_{1093}(2)=364$.  Also
\[
2^{364}\equiv1\pmod{1093^2},\qquad
2^{364}\equiv1064432260\not\equiv1\pmod{1093^3},
\]
so $w_{1093}=2$.  These are the finite checks recorded in [112]; the
two prime-power congruences were recomputed for this manuscript.
They give $D_{1093^4}^{(2)}=1$.
Thus every quadratic weighted binary representation of
$1/1093^4$ is aperiodic.  This is a consequence for one explicitly
checked prime, and makes no assertion about the infinitude of
Wieferich primes.

\begin{corollary}[Rational derivative data and a natural boundary]
\label{disc:naturalboundary}
For binary digits define
\[
 Y=\sum_{n\ge1}\frac{\varepsilon_n}{2^n},\qquad
 F(z)=\sum_{n\ge1}\varepsilon_nz^n.
\]
The following conditions are equivalent: the digits are eventually
periodic; $Y$ is rational; $F$ is a rational function.
If $x$ satisfies Corollary~\ref{disc:cubefree}, then every quadratic
weighted representation of $x$ therefore gives
\[
 Y\notin\mathbb Q,
 \qquad F\text{ has the unit circle as a natural boundary},
 \qquad
 \left.\left(z\frac{d}{dz}\right)^2F(z)\right|_{z=1/2}
 =x\in\mathbb Q.
\]
In particular, $F$ is transcendental over $\mathbb C(z)$.
\end{corollary}

\begin{proof}
The equivalence between eventual periodicity and rationality of $Y$
is the ordinary radix-expansion theorem.  Periodicity makes $F$
rational by the geometric-series formula.  Conversely, a rational
$F$ gives an eventual linear recurrence with constant rational
coefficients.  A state consisting of sufficiently many consecutive
digits determines the next digit.  There are only finitely many such
binary states, so a repeated state forces eventual periodicity.

For a nonperiodic sequence, infinitely many digits are one, and the
integer-coefficient series $F$ has radius of convergence exactly one.
The P\'olya--Carlson theorem \cite{Bell2023} says that such a series is either rational
or has the unit circle as a natural boundary.  The rational case has
been excluded.  An algebraic function has only finitely many possible
singularities and branch points, so cannot have that natural boundary.
The derivative identity follows by termwise differentiation strictly
inside the circle of convergence.
\end{proof}

\begin{corollary}[An exact recurrence criterion for binary sequences]
\label{disc:recurrencecriterion}
For a binary sequence $\varepsilon_n$, the following six conditions
are equivalent:
\begin{enumerate}[label=\textup{(\roman*)}]
\item the sequence is eventually periodic;
\item $\sum_{n\ge1}\varepsilon_n2^{-n}$ is rational;
\item $F(z)=\sum_{n\ge1}\varepsilon_nz^n$ is rational;
\item $F(z)$ is algebraic over $\mathbb C(z)$;
\item $F(z)$ is D-finite over $\mathbb C(z)$;
\item the sequence is P-recursive: for some polynomials
$P_0,\ldots,P_m\in\mathbb C[n]$, not all zero,
\[
 \sum_{j=0}^mP_j(n)\varepsilon_{n+j}=0
\]
for all sufficiently large $n$.
\end{enumerate}
Consequently every quadratic weighted representation in
Corollary~\ref{disc:cubefree} fails all six conditions, including every
nonzero polynomial-coefficient linear recurrence.
\end{corollary}

\begin{proof}
The first three equivalences were proved in
Corollary~\ref{disc:naturalboundary}.  A rational function is algebraic
and D-finite.  An algebraic function cannot have a full circle as a
natural boundary, since its possible branch points and singularities
belong to a finite set.  A D-finite analytic germ satisfies a nonzero
linear differential equation with polynomial coefficients, and can be
continued along paths avoiding the finite zero set of its leading
coefficient.  It therefore cannot have the unit circle as a natural
boundary either.  The P\'olya--Carlson dichotomy in the preceding proof
forces rationality in both cases.

Finally, coefficient extraction converts a polynomial-coefficient
differential equation into an eventual polynomial-coefficient recurrence,
and summing such a recurrence converts it into a polynomial-coefficient
differential equation with at most a polynomial inhomogeneous term.
Applying sufficiently many derivatives removes that term.  This proves
the equivalence of D-finiteness and P-recursiveness used here.
\end{proof}

A further standard consequence links this criterion to corpus [111]:
no such digit generating function can be a diagonal of a rational
function over $\mathbb Q$ in finitely many variables, since these
diagonals are D-finite in characteristic zero.  This is a consequence
of the classical diagonal theorem, not a separate conjecture resolution.

\section{The complete binomial gcd distribution in a prime family}
\label{disc:gcdsection}

Corpus [086] proves a minimum-value criterion by examining the base-prime
digits of three nearby primes.  Its increasing-gap half also yields the
full distribution of the binomial gcd values.

\begin{theorem}[Exact four-value histogram]
\label{disc:histogram}
Let $p<q<r$ be primes and put
\[
 a=q-p,\quad c=r-q,\quad b=a+c=r-p,\quad n=pqr.
\]
Suppose $2b^2<p$ and $c>a$.  Set
\begin{align*}
 A_p&=2(a+b+1)(ab+1)-2,\\
 A_q&=2(c-a)(q-ac+1)-2,\\
 A_r&=(r-b-c+1)(bc+1)-2.
\end{align*}
For $1\le k\le n-1$, the distribution is exactly
\[
\begin{array}{c|c}
\gcd(n,\binom nk)&\text{number of indices}\\\hline
qr&A_p\\
pr&A_q\\
pq&A_r\\
n&n-1-A_p-A_q-A_r.
\end{array}
\]
For $1<k\le n/2$, the first three counts are halved, and the last
count is $(n-3-A_p-A_q-A_r)/2$.
Furthermore, uniformly in this family, all but $O(pb^2)$ of the
interior coefficients are divisible by the entire integer $n$.
The exceptional proportion is $O(b^2/p^2)=O(1/p)$.
\end{theorem}

\begin{proof}
We include the pair-exclusion argument from [086] so the arithmetic
input is explicit.  By symmetry it is enough to consider
$1\le k\le n/2$.  If two distinct primes $s,t\mid n$ fail to divide
$\binom nk$, Lucas' lowest digit forces $k=stj$, where
$1\le j\le(u-1)/2$ and $u$ is the remaining prime.  The necessary
conditions are that $tj$ is digitwise at most $tu$ in base $s$, and
$sj$ is digitwise at most $su$ in base $t$.
The product expansions give
\[
 qr=(1,a+b,ab)_p,\qquad
 pr=(1,c-a-1,q-ac)_q,\qquad
 pq=(r-b-c,bc)_r.
\]
All displayed digits are in range under $2b^2<p$.
We will use $a(a+b),b(a+b),c(b+c)<2b^2<p$.

For the pair $\{p,q\}$, write
$qj=(j+\lfloor aj/p\rfloor)p+(aj\bmod p)$.
Since $j<(p+b)/2$ and $b+2a<p$, its middle digit is less than $p$.
The middle-digit bound gives $j\le a+b$, hence $aj<p$.
The last-digit bound now gives $aj\le ab$, so $j\le b$.
In base $q$ we have $pj=(j-1)q+(q-aj)$.
Its two digit bounds force respectively $j\le c-a$ and $j\ge c$,
a contradiction.

For the pair $\{p,r\}$, similarly
$rj=(j+\lfloor bj/p\rfloor)p+(bj\bmod p)$ has two digits,
since $j<(p+a)/2$ and $a+2b<p$.
The middle-digit bound gives $j\le a+b$, so $bj<p$;
the last-digit bound then yields $bj\le ab$ and $j\le a$.
In base $r$, $pj=(j-1)r+(r-bj)$, and the last-digit bound requires
$r\le b(j+c)\le b^2<p$, again impossible.

For $\{q,r\}$, write
$rj=(j+\lfloor cj/q\rfloor)q+(cj\bmod q)$.
Here $j<p/2$ makes the middle digit less than $q$, so its bound gives
$j\le c-a-1<c$.  Thus $cj<c^2<r$, and in base $r$,
$qj=(j-1)r+(r-cj)$.  Its last-digit bound requires
$r\le c(j+b)<c(b+c)<p$, the final contradiction.
Therefore at least two primes divide every interior coefficient.

For a prime $s$ and base-$s$ digits $n_j$ of $n$, Lucas' theorem gives
\[
 \#\{0\le k\le n:s\nmid\tbinom nk\}=\prod_j(n_j+1).
\]
Indeed each digit of $k$ can independently be chosen from $0$ to
$n_j$.  Excluding $k=0,n$ subtracts two.  The product expansions
give the base-prime digits
\[
 qr=(1,a+b,ab)_p,\qquad
 pr=(1,c-a-1,q-ac)_q,\qquad
 pq=(r-b-c,bc)_r.
\]
Every displayed digit is in range because $2b^2<p$.
Multiplication by the corresponding missing prime appends a zero
digit to obtain $n$.  The Lucas counts are therefore $A_p,A_q,A_r$.
The three nondivisibility sets are disjoint by the first paragraph,
so these are exactly the counts for gcd $n/p,n/q,n/r$; the remaining
indices have gcd $n$.

Since $n$ is odd, the symmetry $k\leftrightarrow n-k$ halves each
interior count.  Excluding $k=1$ removes one only from the gcd-$n$
count.  Finally the three explicit expressions are respectively
$O(b^3)$, $O(pb)$, and $O(pb^2)$.  As $b<p$ and $n\ge p^3$, the
claimed concentration estimates follow.
\end{proof}

For $p=101$, $q=103$, and $r=107$, the hypotheses hold because
$2\cdot6^2=72<101$ and $4>2$.  Here $n=1,113,121$, and for
$2\le k\le556,560$ the four counts are
\[
\begin{array}{r|r}
\gcd(n,\binom nk)&\text{count}\\\hline
10,403&1,224\\
10,807&191\\
11,021&116\\
1,113,121&555,028.
\end{array}
\]
They sum to $556,559$.  An independent computation using Legendre's
factorial valuation formula checked every index in this interval.
This histogram extends the conclusion of [086]; it does not classify
arbitrary squarefree integers or settle the remaining asymptotic
upper-bound question in Erd\H{o}s Problem 700.

\section{Translated reciprocal sums in one number field}
\label{disc:fieldsection}

The rational multilinear-tail argument of corpus [080] admits a number
field extension.  The new input is the norm of a nonzero algebraic
integer; no sign assumption is imposed on its conjugates.

\begin{theorem}[Growth versus common field degree]
\label{disc:fielddegree}
Let $a_1<a_2<\cdots$ be positive integers satisfying
\begin{equation}\label{disc:taildomination}
 \sup_n a_n\sum_{k\ge n}\frac1{a_k}<\infty.
\end{equation}
For rational $h$ avoiding all poles $-a_n$, write
\[
 S(h)=\sum_{n\ge1}\frac1{a_n+h}.
\]
Suppose that, for $d$ distinct rational numbers $h_1,\ldots,h_d$,
all $S(h_i)$ belong to a number field $K$ of degree
$D=[K:\mathbb Q]$.  Then
\[
 \boxed{\log a_n=O\bigl((dD+1)^{n/d}\bigr).}
\]
For integer shifts, after a finite truncation, there is a constant $C$
such that
\[
 \sum_{j=0}^{d-1}\log a_{n+j}
 \le dD\sum_{k<n}\log a_k+C.
\]
\end{theorem}

\begin{proof}
First take integer shifts $h_1<\cdots<h_d$, and discard finitely
many terms so all denominators are positive. The resulting tail values
are denoted by $S(h_i)$ again; each differs from its original value by
a rational number and hence still belongs to $K$. Write
\[
 \alpha_k^{(i)}=a_k+h_i,\quad
 \Pi_n^{(i)}=\prod_{k<n}\alpha_k^{(i)},\quad
 R_n^{(i)}=S(h_i)-\sum_{k<n}\frac1{\alpha_k^{(i)}}.
\]
Choose a positive integer $Q$ for which every $QS(h_i)$ is an
algebraic integer.  Then $Q\Pi_n^{(i)}R_n^{(i)}$ is an algebraic
integer.  Define the strictly positive sum
\[
 O_n=\sum_{\substack{k_1,\ldots,k_d\ge n\\
                    k_1,\ldots,k_d\ \mathrm{distinct}}}
               \prod_{i=1}^d\frac1{\alpha_{k_i}^{(i)}}.
\]

For a block $B\subseteq\{1,\ldots,d\}$ put
$c_{B,i}=\prod_{j\in B\setminus\{i\}}(h_j-h_i)^{-1}$.
Partial fractions give
\[
 \prod_{i\in B}\frac1{x+h_i}
   =\sum_{i\in B}\frac{c_{B,i}}{x+h_i}.
\]
For a set partition $\pi$, let
$\mu(\pi)=\prod_{B\in\pi}(-1)^{|B|-1}(|B|-1)!$.
M\"obius inversion on set partitions, followed by the displayed
partial fractions, gives the exact identity
\begin{equation}\label{disc:moebius}
 O_n=\sum_{\pi}\mu(\pi)
       \prod_{B\in\pi}\left(\sum_{i\in B}c_{B,i}R_n^{(i)}\right).
\end{equation}
The inversion removes precisely the tuples having repeated indices;
all sums converge absolutely.

Let $\Delta=\prod_{i<j}(h_j-h_i)>0$.  In an expanded term of
\eqref{disc:moebius}, the representatives chosen from different blocks
are distinct.  The denominator of the product of their $c_{B,i}$
coefficients is a product over distinct unordered pairs and divides
$\Delta$.  Consequently
\[
 Z_n=M_nO_n,\qquad
 M_n=\Delta Q^d\prod_{i=1}^d\Pi_n^{(i)},
\]
is a nonzero algebraic integer.

For every embedding $\sigma:K\longrightarrow\mathbb C$,
\[
 \sigma(R_n^{(i)})=
 \sigma(S(h_i))-\sum_{k<n}\frac1{\alpha_k^{(i)}}
\]
is uniformly bounded in $n$.  Equation~\eqref{disc:moebius} therefore
gives $|\sigma(O_n)|\le B$ for a fixed $B\ge1$, simultaneously
for all embeddings and all $n$.  The norm of $Z_n$ is a nonzero
rational integer.  At the distinguished real embedding $O_n>0$,
and thus
\[
 1\le\left|\operatorname{Norm}_{K/\mathbb Q}(Z_n)\right|
   =M_n^DO_n\prod_{\sigma\ne\mathrm{id}}|\sigma(O_n)|
   \le M_n^DB^{D-1}O_n.
\]

Put $P_n=\sum_{k<n}\log a_k$.  The convergence of $\sum1/a_k$
implies that the partial products for every fixed shift differ from
$\prod_{k<n}a_k$ by uniformly bounded positive factors.  Hence
$M_n\le C_1e^{dP_n}$.  Symmetrization and then removal of the
ordering restrictions on distinct indices give
\[
 O_n\le d!\prod_{j=0}^{d-1}
             \left(\sum_{k\ge n+j}\frac1{a_k+h_1}\right)
       \le C_2\prod_{j=0}^{d-1}a_{n+j}^{-1},
\]
where the last inequality is \eqref{disc:taildomination}.
Combining the estimates and taking logarithms proves
\[
 P_{n+d}\le(dD+1)P_n+C.
\]
Iteration on each residue class modulo $d$ gives the required bound.

For rational shifts take a common denominator $w$, write $h_i=u_i/w$,
and replace $a_n$ by $b_n=wa_n$.  The translated sums become
$S(h_i)/w\in K$, condition~\eqref{disc:taildomination} is unchanged,
and $\log b_n=\log a_n+O(1)$.  The integer-shift result applies.
\end{proof}

\begin{corollary}[Finitely many translated values in each fixed field]
\label{disc:fieldcount}
Under \eqref{disc:taildomination}, suppose
\[
 \lambda=\limsup_{n\to\infty}(\log a_n)^{1/n}>1.
\]
If $d$ distinct translated values $S(h_i)$ lie in a number field of
degree $D$, then
\[
 \lambda^d\le dD+1.
\]
Thus each fixed number field contains only finitely many values
$S(h)$ at rational shifts avoiding the poles. If $1<\lambda<\infty$
and $d$ selected values are all algebraic, then
\[
 \boxed{
 [\mathbb Q(S(h_1),\ldots,S(h_d)):\mathbb Q]
\ge\left\lceil\frac{\lambda^d-1}{d}\right\rceil.}
\]
If $\lambda=\infty$, every allowed rational-shift value $S(h)$ is
transcendental.
\end{corollary}

\begin{proof}
The growth bound of Theorem~\ref{disc:fielddegree} gives
$\lambda\le(dD+1)^{1/d}$.  For fixed $D$, this fails for all
sufficiently large $d$, proving finiteness.  Apply the same inequality
to the field generated by the selected algebraic values.
When $\lambda=\infty$, apply the theorem with $d=1$ to the number
field generated by any putatively algebraic value; the resulting
finite bound $\lambda\le D+1$ is impossible.
\end{proof}

The conclusion applies uniformly to every positive integer sequence
with $a_n/2^{2^n}\to1$, the perturbation class in Erd\H{o}s
Problem 263 \cite{Erdos263}. Indeed $a_{n+1}/a_n\to\infty$, so
eventual monotonicity and the tail hypothesis follow after discarding
a finite prefix, while $\log a_n=2^n\log2+o(1)$ gives $\lambda=2$.
If all values at $d$ distinct rational shifts are algebraic, their
compositum has degree at least $\lceil(2^d-1)/d\rceil$.
For $d=5,10,16$ these lower bounds are $7,103,4096$, respectively.
At most one allowed shift value belongs to $\mathbb Q$, and at most
two belong to any specified quadratic field. In particular no three
distinct rational shifts can all give values in the same quadratic
field, for any sequence in this entire perturbation class.

An explicit perturbed example is $a_n=2^{2^n}+n$. It has
$a_{n+1}\ge3a_n$ and hence tail constant at most $3/2$.
These are common-field restrictions; they do not determine the
arithmetic nature of a specified $S(0)$ throughout this class or
settle Erd\H{o}s Problems 263 or 265. We deliberately do not label
the unperturbed Mahler/Fermat special values as newly unresolved:
classical transcendence results already apply to important instances
of those values \cite{Coons2011}.

For comparison, $a_n=2^{n!}$ satisfies the tail hypothesis and has
$\lambda=\infty$. Thus every number
$\sum_{n\ge1}(2^{n!}+h)^{-1}$ at an allowed rational shift $h$ is
transcendental. This is the classical rapid-growth mechanism; the
fixed-field degree restriction above reaches the slower sequence
$2^{2^n}$ where that one-variable conclusion does not follow.

\paragraph{Checks.}
Exact rational arithmetic verified the quadratic valuation identity
in $4,330$ prime-valuation cases coming from primitive binary periods at
most eight and prefix lengths from zero through four.  The binomial
histogram example was independently verified at every relevant index
using Legendre's formula, rather than the Lucas count used in the
proof.  These checks supplement the all-index arguments above.


\section{A geometric converse to polynomial primitive towers}
\label{sec:geometry}

\paragraph{Source and additional step.}
Corpus [113] proves that polynomial covers possess algebraic primitive
towers, and [114] organizes stability by homology and monodromy.  We
extract a converse that does not need the representation-theoretic
classification: universal integrability detects the genus and the
number of points above infinity.  Its proof uses trace descent,
approximation of finite jets, and the classical de Rham theorem for
affine curves.  These are structural consequences of classical tools;
no assertion of historical priority is made.

Let $X$ be a smooth connected projective complex curve, let
$x\in K=\mathbb C(X)$ be nonconstant, and let $S$ be the nonempty set
of poles of $x$. Put $U=X\setminus S$, $A=\mathcal O(U)$, and
$D=d/dx$. Write $g$ for the genus of $X$ and $s=|S|$.
All assertions concern characteristic zero. An algebraic primitive is
allowed initially in a finite algebraic extension of $K$; elementary
primitives are allowed in elementary extensions with the same constants.

\begin{lemma}[Trace descent and absence of new finite poles]
\label{geom:trace}
If $f\in K$ has an algebraic $n$th primitive, it has one in $K$.
If $f\in A$ and $Dh=f$ for $h\in K$, then $h\in A$.
\end{lemma}
\begin{proof}
For a finite extension $L/K$, the derivation extends uniquely and
commutes with all $K$-embeddings. Consequently it commutes with trace.
If $D^n h=f$, then $D^n(\operatorname{Tr}_{L/K}(h)/[L:K])=f$.
For the second assertion, $dh=f\,dx$ has no pole on $U$. A pole of
order $m>0$ of $h$ would give a pole of order $m+1$ of $dh$, which
is impossible.
\end{proof}

\begin{theorem}[Exact dimension of the algebraic-integration obstruction]
\label{geom:codim}
Let
\[
 I_x=\{f\in A: f\text{ has an algebraic primitive with respect to }x\}.
\]
Then
\[
 \boxed{\quad \dim_{\mathbb C}(A/I_x)=2g+s-1.\quad}
\]
More precisely, $f\mapsto[f\,dx]$ induces an isomorphism
$A/I_x\simeq H^1_{\mathrm{dR}}(U)$.
\end{theorem}
\begin{proof}
Because $U$ is a smooth affine curve,
$H^1_{\mathrm{dR}}(U)=\Omega^1_A/dA$. The kernel of the indicated map
is $I_x$, by Lemma~\ref{geom:trace}.
To prove surjectivity, take $\omega\in\Omega^1_A$.
The zeros of $dx$ on $U$ form a finite set. At such a point $P$, choose
a local parameter $z$ and put $m_P=\operatorname{ord}_P(dx)\ge1$.
Choose a truncated Taylor primitive $h_P$ of $\omega$ through degree
$m_P$, so that $\omega-dh_P$ vanishes to order at least $m_P$.
The Chinese remainder theorem in $A$ supplies an $h\in A$ with
$h-h_P\in\mathfrak m_P^{m_P+1}$ at every chosen point.
Therefore $(\omega-dh)/dx$ has no pole anywhere on $U$, and equals
some $f\in A$. Thus $[\omega]=[f\,dx]$.
The classical algebraic de Rham comparison and the topology of a
genus-$g$ surface with $s$ punctures give
$\dim H^1_{\mathrm{dR}}(U)=2g+s-1$.
\end{proof}

\begin{theorem}[Universal algebraic stability characterizes polynomial covers]
\label{geom:polynomial}
The following statements are equivalent:
\begin{enumerate}
\item Every $f\in A$ has an algebraic primitive.
\item Every $f\in A$ has a coherent infinite tower of algebraic primitives.
\item $g=0$ and $s=1$.
\item There is a coordinate $t$ with $K=\mathbb C(t)$ and
      $x=P(t)$ for a nonconstant polynomial $P\in\mathbb C[t]$.
\end{enumerate}
In the affirmative case every tower can be chosen in $K$.
\end{theorem}
\begin{proof}
The second assertion implies the first. Theorem~\ref{geom:codim}
makes the first equivalent to $2g+s-1=0$, hence to $g=0,s=1$.
A genus-zero smooth projective curve over $\mathbb C$ is the projective
line. Choosing its unique point in $S$ as infinity gives
$A=\mathbb C[t]$ and makes $x$ a polynomial in $t$; the converse is
immediate. Finally, on a polynomial cover define
\[
 J_PQ(t)=\int_0^t P'(u)Q(u)\,du,\qquad Q\in\mathbb C[t].
\]
This is polynomial and $D(J_PQ)=Q$. Hence $J_P^nQ$ gives a coherent
infinite tower for every $Q\in A$.
\end{proof}

\begin{theorem}[Universal elementary stability characterizes genus zero]
\label{geom:elementary}
The following are equivalent:
\begin{enumerate}
\item Every $f\in K$ has an elementary primitive with respect to $x$.
\item Every $f\in K$ has elementary primitives of every order.
\item $g=0$.
\end{enumerate}
When $g>0$, the $g$-dimensional space
$\{\omega/dx:\omega\in H^0(X,\Omega_X^1)\}$ consists, except for zero,
entirely of obstructions already at the first integration.
\end{theorem}
\begin{proof}
If $g=0$, write $K=\mathbb C(t)$ and $x=R(t)$ with $R$ rational.
For every $j\ge0$, the differential $x^j f\,dx$ is rational in $t$,
and hence has an elementary primitive $H_j$ by partial fractions.
For $n\ge1$, repeated integration by parts gives
\[
 F_n=\frac1{(n-1)!}\sum_{j=0}^{n-1}(-1)^j\binom{n-1}{j}
                 x^{n-1-j}H_j,
 \qquad D^nF_n=f.
\]
Using the same chosen primitives $H_j$ for all $n$ gives
$DF_n=F_{n-1}$ for $n\ge2$ directly, so these form a coherent tower.

Conversely, suppose $g>0$ and choose a nonzero holomorphic differential
$\omega$. Put $f=\omega/dx\in K$. Liouville's theorem for integration
of algebraic functions, followed by trace and norm from its finite
algebraic extension, would express an elementary-integrable $\omega$ as
\[
 \omega=dh+\sum_{i=1}^r c_i\,du_i/u_i,
 \qquad h\in K,\quad u_i\in K^*,\quad c_i\in\mathbb C.
\]
A pole of $h$ would produce a pole of order at least two in $dh$,
which the logarithmic differentials cannot cancel. Thus $h$ is constant.
The vanishing residues of $\omega$ then give
$\sum_i c_i\operatorname{ord}_P(u_i)=0$ for all $P$.
The finite matrix of these valuations has integer entries, so its
complex kernel is generated by rational vectors. Multiplying a
rational kernel vector by a common denominator yields integers $a_i$
with $\operatorname{div}(\prod_i u_i^{a_i})=0$. The product is constant
on $X$, and hence $\sum_i a_i\,du_i/u_i=0$. Every residue-free
linear combination in the display is therefore zero, contrary to
$\omega\ne0$. This proves both necessity and the last statement.
\end{proof}

\paragraph{An explicit elliptic obstruction.}
For $y^2=x^3-x$, the smooth projective curve has genus one and one
point above infinity. Therefore $A=\mathbb C[x,y]/(y^2-x^3+x)$ has
exactly two independent one-step algebraic-integration obstructions.
The identities
\begin{align*}
 d\bigl((3x^2-2)y\bigr)&=dx/y+(21/2)xy\,dx,\\
 d(xy)&=x\,dx/y+(5/2)y\,dx
\end{align*}
give concrete representatives. Since $dx/y$ is a nonzero holomorphic
differential, $xy$ has no elementary primitive. The classes of $dx/y$
and $x\,dx/y$ are independent: if the
second were a multiple of the first plus an exact differential, its
primitive would have a unique simple pole at infinity, which is
impossible on a genus-one curve by Riemann--Roch. Thus $xy$ and $y$
represent the two independent algebraic obstructions. For
$y^2=P(x)$ with $P$ squarefree of degree $d\ge2$, the dimension in
Theorem~\ref{geom:codim} is $d-1$. This count is about first primitives;
it must not be confused with the full repeated-integration
classification of [114].

\paragraph{What this connection establishes.}
Universal integrability on an entire algebraic function ring or field
is equivalent to a precise geometric condition. It does not prove
elementary stability of arbitrary transcendental elementary inputs,
and it does not certify the broader claims of corpus [119]--[121].


\section{A positivity threshold, finite moment tests, and a martingale comparison}
\label{spec:section}

This section develops the inverse-Mellin argument of corpus manuscript
[116].  The link between complete monotonicity and a Stieltjes moment
sequence is already explicit in Simon's paper \cite[Section~3]{Simon2021}.
The additional conclusions here concern the structure of the actual
positivity boundary, its finite moment approximations, a parameter
evolution equation, and a probabilistic comparison.  They do not provide
a closed formula for that boundary.  The strict separation from the
logarithmic infinite-divisibility boundary is the starting result of
[116]; the analytic ingredients needed for it are included below.

Fix throughout
\[
 \gamma>1,\qquad 0<\alpha<1/\gamma,\qquad \nu=1-\alpha\gamma\in(0,1).
\]
For \(\beta>0\), write
\[
 F_\beta(z)=\sum_{n=0}^{\infty}
             \frac{z^n}{\Gamma(\beta+\alpha n)^\gamma},
 \qquad
 m_n(\beta)=\frac{n!\,\Gamma(\beta)^\gamma}
                   {\Gamma(\beta+\alpha n)^\gamma}.
\]
The power series for \(F_\beta\) is entire.  Set
\begin{equation}\label{spec:ID-boundary}
 \beta_{\rm ID}=\max_{t>0}
 \left[\alpha+\frac{\alpha}{t}
  \log\!\left(\frac{\gamma(1-e^{-t})}{1-e^{-t/\alpha}}\right)\right].
\end{equation}
This is the threshold for the scalar inequality in Simon's conjecture,
and is not the threshold for complete monotonicity.

\subsection{Analytic preparation}

For \(\beta\ge\alpha\) and \(\Re s>-1\), use the holomorphic logarithm
of Gamma in the right half-plane to define
\[
 M_\beta(s)=\Gamma(1+s)\Gamma(\beta)^\gamma
       \exp\{-\gamma\operatorname{Log}\Gamma(\beta+\alpha s)\}.
\]
The following real, initially signed, density is well defined:
\begin{equation}\label{spec:q}
 q_\beta(x)=\frac{1}{2\pi i x}
       \int_{-i\infty}^{i\infty}M_\beta(s)x^{-s}\,ds,
 \qquad x>0.
\end{equation}

\begin{lemma}[Uniform endpoint estimates]\label{spec:endpoints}
The function \(q_\beta(x)\) is jointly real analytic in
\(\beta>\alpha\) and \(x>0\).  It has total signed mass one and
\begin{equation}\label{spec:mellin}
 \int_0^\infty x^s q_\beta(x)\,dx=M_\beta(s),\qquad \Re s>-1.
\end{equation}
For every compact \(K\subset(\alpha,\infty)\), uniformly for \(\beta\in K\),
\begin{align}
 q_\beta(x)&=\frac{\Gamma(\beta)^\gamma}
                   {\Gamma(\beta-\alpha)^\gamma}+O_K(x^\eta),
                   &&x\downarrow0,\label{spec:zero-tail}\\
 q_\beta(x)&=\frac{C_\beta}{x\sqrt{2\pi\nu}}
           r^{\kappa_\beta+1/2}e^{-\nu r}(1+o_K(1)),
                   &&x\to\infty,\label{spec:infty-tail}
\end{align}
where
\[
 \kappa_\beta=\frac{1+\gamma}{2}-\gamma\beta,
 \qquad r=(\alpha^{\alpha\gamma}x)^{1/\nu},
\]
\[
 C_\beta=(2\pi)^{(1-\gamma)/2}\Gamma(\beta)^\gamma
                       \alpha^{\gamma(1/2-\beta)},
\]
where any \(0<\eta<\min(1,\min K/\alpha-1)\) is allowed.
In particular, the densities are positive near both endpoints uniformly
on \(K\), and their absolute values have every positive exponential moment.
At \(\beta=\alpha\), the same Mellin identity holds, the upper-end estimate
continues to hold, and \(q_\alpha(x)=O_\varepsilon(x^{-\varepsilon})\)
as \(x\downarrow0\), for every \(\varepsilon>0\).
\end{lemma}

\begin{proof}
Vertical Stirling estimates give a polynomial in \(|\Im s|\) times
\(e^{-\pi\nu|\Im s|/2}\) on each permissible fixed vertical line.
This proves absolute convergence, differentiation under the integral,
and analyticity in the logarithmic variable \(\log x\); small complex
neighborhoods of \(\beta\) have the same estimates.  Moving the contour
to \(\Re s=-1-\eta\) crosses only the numerator pole at \(s=-1\) and
gives \eqref{spec:zero-tail}.  When \(\beta=\alpha\), moving instead to
\(\Re s=-1+\varepsilon\) gives the stated bound without crossing a pole.

For completeness the uniform large-\(x\) step is as follows.  Stirling's
formula on the right half-plane gives
\[
 M_\beta(s)=C_\beta s^{\kappa_\beta}
 \exp\{\nu s(\log s-1)-\alpha\gamma s\log\alpha\}
 (1+O_K(|s|^{-1})).
\]
Shift to \(\Re s=r\), where
\(\log x+\alpha\gamma\log\alpha=\nu\log r\), and set
\(s=r+i\sqrt r\,u\).  The normalized integrand tends to
\(e^{-\nu u^2/2}\).  If \(v=u/\sqrt r\), its modulus is bounded by
\[
 D_K(1+v^2)^{\kappa_\beta/2}e^{\nu r H(v)},\qquad
 H(v)=\tfrac12\log(1+v^2)-v\arctan v.
\]
The bounds \(H(v)\le-v^2/4\) for \(|v|\le1\) and
\(H(v)\le H(1)-\pi(|v|-1)/4\) for \(|v|\ge1\) prove uniform
dominated convergence on the central part and exponential negligibility
of the outer part.  This proves \eqref{spec:infty-tail}.  These are the
uniform saddle estimates used in [116], based on \cite[Section~5.11]{DLMF}.

The endpoint bounds imply integrability of the signed density and of
\(x^s q_\beta(x)\) for \(\Re s>-1\).  In the coordinate \(y=\log x\),
Fourier inversion identifies the transform of \(e^yq_\beta(e^y)\)
with \(M_\beta(iu)\).  Analytic continuation in \(s\) proves
\eqref{spec:mellin}, including mass one at \(s=0\).
Since \(1/\nu>1\), the upper tail also gives all positive exponential
moments of \(|q_\beta|\).
\end{proof}

For \(\beta\ge\alpha\), absolute integration of the exponential series
therefore gives
\begin{equation}\label{spec:signed-laplace}
 \int_0^\infty e^{zx}q_\beta(x)\,dx
       =\Gamma(\beta)^\gamma F_\beta(z),\qquad z\in\mathbb C.
\end{equation}
Bernstein's theorem and uniqueness of Laplace transforms of finite
signed measures imply the exact equivalence
\begin{equation}\label{spec:CM-positive}
 F_\beta(-x)\text{ is completely monotone}
 \quad\Longleftrightarrow\quad q_\beta(x)\ge0\quad(x>0).
\end{equation}

\subsection{A parameter coupling and a single positivity boundary}

\begin{lemma}[Beta factors]\label{spec:beta-factor}
For every \(\beta,\delta>0\), there is a random variable
\(Z_{\beta,\delta}\) with \(0<Z_{\beta,\delta}<1\) such that
\begin{equation}\label{spec:Z-mellin}
 \mathbb E Z_{\beta,\delta}^{\,s}
 =\left[
   \frac{\Gamma(\beta+\delta)\Gamma(\beta+\alpha s)}
        {\Gamma(\beta)\Gamma(\beta+\delta+\alpha s)}
   \right]^\gamma,\qquad \Re s\ge0.
\end{equation}
Consequently the complete-monotonicity property, once true at \(\beta\),
is true at every larger parameter.  When both parameters belong to
the positive-density region, they can be coupled by
\begin{equation}\label{spec:coupling}
 X_{\beta+\delta}\overset{d}=X_\beta Z_{\beta,\delta},
 \qquad X_\beta\ \hbox{independent of }Z_{\beta,\delta}.
\end{equation}
\end{lemma}

\begin{proof}
The logarithm of the right side of \eqref{spec:Z-mellin} is
\[
 \gamma\int_0^\infty (e^{-\alpha st}-1)
 \frac{e^{-\beta t}(1-e^{-\delta t})}{t(1-e^{-t})}\,dt.
\]
Its positive kernel is a subordinator Levy measure: near zero it is
\(\gamma\delta/t+O(1)\), and it decays exponentially at infinity.
Exponentiating the negative of that finite-valued subordinator defines
\(Z\).  Equivalently, it is the multiplicative convolution power
\(\gamma\) of \(B^\alpha\), where \(B\) has the beta law
\(\operatorname{Beta}(\beta,\delta)\).  Its law has no atom at one.
The gamma-ratio identity gives
\(m_n(\beta+\delta)=m_n(\beta)\mathbb E Z^n\).

If complete monotonicity holds at \(\beta\), Bernstein's law has
moments \(m_n(\beta)\).  The product with an independent \(Z\) has
moments \(m_n(\beta+\delta)\).  The entire series for \(F_{\beta+\delta}\)
shows by monotone convergence that this product has all positive
exponential moments and moment generating function
\(\Gamma(\beta+\delta)^\gamma F_{\beta+\delta}\).  Thus its Laplace
transform proves complete monotonicity at \(\beta+\delta\).
Moment or Laplace-transform uniqueness also proves \eqref{spec:coupling}.
\end{proof}

\begin{theorem}[One true threshold, with strict interior positivity]
\label{spec:threshold}
There is a unique number \(\beta_c=\beta_c(\alpha,\gamma)\) such that
\begin{equation}\label{spec:threshold-range}
 \alpha<\beta_c<\beta_{\rm ID}<\frac{\alpha(1+\gamma)}2,
\end{equation}
and
\[
 F_\beta(-x)\text{ is completely monotone}
       \quad\Longleftrightarrow\quad \beta\in[\beta_c,\infty).
\]
The density is strictly positive at every \(x>0\) whenever
\(\beta>\beta_c\).  At \(\beta_c\) it is nonnegative and has a
nonempty finite set of zeros in \((0,\infty)\); all these zeros
have finite even order.
\end{theorem}

\begin{proof}
At \(\beta=\alpha\), if \(q_\alpha\) were nonnegative, then
\eqref{spec:mellin} and Fatou's lemma would give
\[
 0<\int_0^\infty x^{-1}q_\alpha(x)\,dx
 \le\liminf_{s\downarrow-1}M_\alpha(s)=0,
\]
because \(M_\alpha(s)\) is a positive constant times
\((1+s)^{\gamma-1}(1+o(1))\).  The left integral is allowed to be
infinite, and is strictly positive since the mass is one.  Hence
complete monotonicity fails at \(\alpha\).  By the beta-factor lemma
it also fails below \(\alpha\).

We recall why the upper endpoint in \eqref{spec:threshold-range} lies
strictly above some positive-density parameters.  Put
\[
 \phi_\beta(t)=\frac{e^{-t}}{1-e^{-t}}
       -\gamma\frac{e^{-\beta t/\alpha}}{1-e^{-t/\alpha}}.
\]
The condition \(\phi_\beta\ge0\) is equivalent to
\(\beta\ge\beta_{\rm ID}\).  The maximand in
\eqref{spec:ID-boundary} tends to \(-\infty\) at zero and to
\(\alpha\) from above at infinity; its maximum is attained.
It exceeds \(\alpha\).  At \(\beta_U=\alpha(1+\gamma)/2\),
\[
 \phi_{\beta_U}(t)>
 \frac{e^{-t/2}}2\left[
  \frac1{\sinh(t/2)}-\frac\gamma{\sinh(\gamma t/2)}\right]>0.
\]
The attained maximum is consequently smaller than \(\beta_U\).

At \(\beta=\beta_{\rm ID}\), the logarithmic gamma identity
\[
 \log M_\beta(s)
  =\bigl(\psi(1)-\alpha\gamma\psi(\beta)\bigr)s
   +\int_0^\infty(e^{-st}-1+st)\phi_\beta(t)\,\frac{dt}{t}
\]
is a Levy--Khintchine representation.  Thus \(M_\beta(iu)\) is
an infinitely divisible characteristic function with exponential
Fourier decay.  A convolution square root has an analytic nonnegative
density on the real line.  It is not identically zero, so its real
zeros are isolated; its convolution with itself is strictly positive
everywhere.  Fourier uniqueness then proves \(q_\beta(x)>0\) at all
positive \(x\).  Lemma~\ref{spec:endpoints} supplies uniform positive
tails in a compact parameter neighborhood.  The minimum on the
remaining compact \(x\)-interval stays positive for all sufficiently
nearby parameters.  This recovers the strict-gap argument of [116]:
some entire interval below \(\beta_{\rm ID}\) is admissible.

The admissible set is an upper set by Lemma~\ref{spec:beta-factor} and
is closed in \((0,\infty)\), since every derivative of \(F_\beta(-x)\)
depends continuously on \(\beta\).  It is therefore
\([\beta_c,\infty)\), with \(\beta_c>\alpha\); otherwise closedness
would make \(\alpha\) admissible.  The strict gap gives
\(\beta_c<\beta_{\rm ID}\).

To obtain strict positivity above the boundary, split the factor in
\eqref{spec:Z-mellin} as \(Z=B^\alpha W\), where
\(B\sim\operatorname{Beta}(\beta_c,\delta)\), and \(W>0\) is the
independent multiplicative convolution power \(\gamma-1\).
The density \(q_{\beta_c}\) is analytic and nonnegative, with mass
one, so its zeros are isolated.  For every \(x>0\),
\[
 q_{\beta_c+\delta}(x)
   =\mathbb E\bigl[Z^{-1}q_{\beta_c}(x/Z)\bigr]>0.
\]
Indeed, conditional on \(W\), the continuous beta variable avoids
the discrete zero set with probability one.  The integrand is bounded
for fixed \(x\) by the upper-tail estimate, so the formula is finite.

If \(q_{\beta_c}\) had no zero, uniform positive tails and compactness
would preserve positivity below \(\beta_c\), a contradiction.  All its
zeros lie in a compact subinterval by Lemma~\ref{spec:endpoints};
analyticity makes their number finite, and nonnegativity makes each
order even with positive leading coefficient.
\end{proof}

\subsection{The boundary is crossed transversally in the parameter}

\begin{theorem}[Evolution equation and local loss of positivity]
\label{spec:contact}
For \(\beta>\alpha\),
\begin{equation}\label{spec:generator}
 \partial_\beta q_\beta(x)
 =\gamma\int_0^\infty\frac{e^{-\beta t}}{1-e^{-t}}
       \bigl[e^{\alpha t}q_\beta(xe^{\alpha t})-q_\beta(x)\bigr]\,dt.
\end{equation}
At every zero \(x_*\) of \(q_{\beta_c}\),
\[
 q_{\beta_c}(x_*)=\partial_xq_{\beta_c}(x_*)=0,
 \qquad b:=\partial_\beta q_{\beta_c}(x_*)>0.
\]
If this zero has order \(2m\), and
\(a=\partial_x^{2m}q_{\beta_c}(x_*)/(2m)!>0\), then as
\(\varepsilon\downarrow0\) a single negative interval is born near
\(x_*\) for \(q_{\beta_c-\varepsilon}\).  Its two endpoints satisfy
\begin{equation}\label{spec:splitting}
 x_\pm(\varepsilon)
  =x_*\pm(b/a)^{1/(2m)}\varepsilon^{1/(2m)}
       +o(\varepsilon^{1/(2m)}).
\end{equation}
Its negative mass is
\begin{equation}\label{spec:negative-mass}
 \int_{x_-(\varepsilon)}^{x_+(\varepsilon)}
       -q_{\beta_c-\varepsilon}(x)\,dx
 =\frac{4m}{2m+1}\,b(b/a)^{1/(2m)}
       \varepsilon^{1+1/(2m)}(1+o(1)).
\end{equation}
In particular, an ordinary quadratic contact has a square-root
interval width and negative mass of order \(\varepsilon^{3/2}\).
No assertion that all contacts are quadratic is needed.
\end{theorem}

\begin{proof}
The digamma integral gives
\[
 \partial_\beta\log M_\beta(s)
  =\gamma\int_0^\infty
       \frac{e^{-\beta t}(e^{-\alpha st}-1)}{1-e^{-t}}\,dt.
\]
Mellin inversion gives \eqref{spec:generator}.  Interchanging integrals
is justified near \(t=0\) by the cancellation in the brackets and
integrability of \((1+|u|)|M_\beta(iu)|\); at infinity the exponential
kernel suffices.  The resulting real-space integral is also absolutely
convergent, by analyticity at small \(t\) and the upper tail at large
\(t\).

At a zero of the critical density the subtracted term vanishes.
The remaining integrand is nonnegative and is positive except at a
discrete set of \(t>0\).  It has positive integral.  Joint analyticity
therefore gives, with \(h=x-x_*\),
\[
 q_{\beta_c-\varepsilon}(x_*+h)
    =a h^{2m}-b\varepsilon
        +O(h^{2m+1}+\varepsilon |h|+\varepsilon^2).
\]
After \(h=\varepsilon^{1/(2m)}u\) and division by \(\varepsilon\),
the local function tends to \(au^{2m}-b\), together with its derivatives
on compact sets.  Analytic root separation, or the Weierstrass preparation
theorem, gives exactly two nearby real roots, with the asymptotics
\eqref{spec:splitting}; the other local roots are nonreal.  The sign
between the real roots is negative.  Integration of the scaled limit
over \([-(b/a)^{1/(2m)},(b/a)^{1/(2m)}]\) gives
\eqref{spec:negative-mass}.
\end{proof}

\begin{corollary}[The largest-contact characterization]\label{spec:largest-contact}
Define
\[
 \mathcal T=\{\beta\in(\alpha,\beta_{\rm ID}):
       q_\beta(x)=\partial_xq_\beta(x)=0
       \text{ for some }x>0\}.
\]
This set is nonempty, has a largest element, and
\[
              \max\mathcal T=\beta_c.
\]
Consequently complete monotonicity is equivalent to
\(\beta\ge\max\mathcal T\).  The original infinite collection of
derivative inequalities can thus be replaced by the largest-parameter
solution of the explicit two-equation contact system.  This is an exact
implicit characterization, not an evaluation of its largest solution.
\end{corollary}

\begin{proof}
Theorem~\ref{spec:threshold} supplies a contact at \(\beta_c\).
Strict positivity at every larger parameter excludes all larger
solutions of the contact system.  Smaller contacts, if any, do not
affect its largest parameter.
\end{proof}

\subsection{An exact hierarchy of finite moment obstructions}

For \(N\ge1\), put
\[
 H_N(\beta)=(m_{i+j}(\beta))_{i,j=0}^N,
 \qquad K_N(\beta)=(m_{i+j+1}(\beta))_{i,j=0}^N.
\]

\begin{theorem}[Finite thresholds converge, but never reach the boundary]
\label{spec:hankel}
For each \(N\ge1\) there is a finite \(b_N>0\) such that
\[
 H_N(\beta)\succeq0\ \hbox{and}\ K_N(\beta)\succeq0
       \quad\Longleftrightarrow\quad \beta\ge b_N.
\]
These thresholds satisfy
\begin{equation}\label{spec:hankel-limit}
 b_N\le b_{N+1}<\beta_c,
 \qquad \lim_{N\to\infty}b_N=\beta_c.
\end{equation}
Thus every non-CM parameter has a finite negative quadratic-form
witness in a Hankel or shifted Hankel matrix.  Nevertheless every
fixed finite collection of these tests accepts some parameters below
the true boundary.
\end{theorem}

\begin{proof}
The beta-factor identity makes \(H_N(\beta+\delta)\) the entrywise
product of \(H_N(\beta)\) with the moment matrix of \(Z_{\beta,\delta}\).
The analogous statement holds for \(K_N\).  Schur's product theorem
therefore makes the simultaneous positive-semidefinite set an upper
set.  It is closed, and is nonempty by Theorem~\ref{spec:threshold}.
As \(\beta\downarrow0\), \(m_1\) and \(m_2\) are positive constants
times \(\beta^{-\gamma}(1+o(1))\), so
\(m_2-m_1^2<0\).  Hence its left endpoint \(b_N\) is positive.

At \(\beta_c\), both matrices are strictly positive definite:
their quadratic forms are respectively
\(\int p(x)^2q_{\beta_c}(x)\,dx\) and
\(\int xp(x)^2q_{\beta_c}(x)\,dx\), which are positive for every
nonzero polynomial \(p\).  Continuity of finitely many eigenvalues
preserves this property in a parameter neighborhood, proving
\(b_N<\beta_c\).  Taking principal submatrices gives monotonicity in
\(N\).

Suppose \(b_*:=\sup_N b_N<\beta_c\), and choose
\(b_*<\beta<\beta_c\).  Every finite Hankel and shifted Hankel
matrix is then positive semidefinite.  The classical Stieltjes moment
theorem gives a probability measure on \([0,\infty)\) with moments
\(m_n(\beta)\).  For every \(t>0\), monotone convergence gives
\[
 \mathbb E e^{tX}
   =\sum_{n=0}^\infty\frac{m_n(\beta)t^n}{n!}
   =\Gamma(\beta)^\gamma F_\beta(t)<\infty.
\]
Absolute expansion of \(e^{-tX}\) therefore identifies its Laplace
transform with \(\Gamma(\beta)^\gamma F_\beta(-t)\), contradicting
the definition of \(\beta_c\).  This proves \eqref{spec:hankel-limit}.
The moment characterization itself is classical and already noted
in \cite[Section~3]{Simon2021}; the strict finite-threshold behavior
uses the parameter coupling and the analytic critical density.
\end{proof}

The first unshifted test gives the explicit necessary condition
\begin{equation}\label{spec:variance-bound}
 \frac{\Gamma(\beta)\Gamma(\beta+2\alpha)}
      {\Gamma(\beta+\alpha)^2}\le2^{1/\gamma}.
\end{equation}
The left side strictly decreases from infinity to one as \(\beta\)
runs from zero to infinity: differentiating its logarithm gives
\(\psi(\beta)+\psi(\beta+2\alpha)-2\psi(\beta+\alpha)<0\), by
strict concavity of the digamma function.  Thus this single necessary
condition already has a unique, explicit one-dimensional threshold.
The shifted \(2\times2\) test similarly requires
\[
 \frac{\Gamma(\beta+\alpha)\Gamma(\beta+3\alpha)}
      {\Gamma(\beta+2\alpha)^2}\le(3/2)^{1/\gamma}.
\]

\subsection{A concrete rational obstruction}

\begin{theorem}[Certified lower bound at \(\alpha=1/3,\gamma=2\)]
\label{spec:rational-witness}
The actual complete-monotonicity threshold satisfies
\[
       \beta_c(1/3,2)>
            \frac{428292245841}{10^{12}}=0.428292245841.
\]
Two explicit rational negative-density witnesses are certified.  The
simpler one is
\begin{equation}\label{spec:q-negative}
 -\frac1{2000}<q_{107/250}(1/2)<-\frac1{2500}<0.
\end{equation}
The sharper one, with
\(\beta_*=428292245841/10^{12}\) and
\(x_*=511574841194/10^{12}\), is
\[
                 -10^{-12}<q_{\beta_*}(x_*)<-10^{-13}<0.
\]
The bound is certified by a finite interval computation and the
explicit analytic remainder estimate below.
\end{theorem}

\begin{proof}
For \(\alpha=1/3\) and \(\gamma=2\), the density has the convergent
entire series
\begin{equation}\label{spec:q-series}
 q_\beta(x)=\Gamma(\beta)^2\sum_{n=0}^\infty
    \frac{(-x)^n}{n!\,\Gamma(\beta-(n+1)/3)^2}.
\end{equation}
Here a reciprocal Gamma value at a nonpositive integer is interpreted
as zero.  One justification, avoiding a formal closure of an infinite
Mellin contour, is to shift the contour finitely many times to read
off \(q(0),q'(0),q''(0)\), and then use the Gamma shift relation in
the Mellin transform.  It gives the ordinary differential equation
\[
 9q'''(x)+x^2q''(x)+(9-6\beta)xq'(x)+(4-3\beta)^2q(x)=0.
\]
The finite contour shifts give smooth extension at zero and the
initial values in \eqref{spec:q-series}.  The equation is regular
at zero and has an analytic solution uniquely determined by these
three values.  Its coefficient recurrence is exactly the recurrence
in \eqref{spec:q-series}, whose three residue-class series are entire
by the ratio test.  Uniqueness proves \eqref{spec:q-series} globally.

Put \(\beta=107/250\), \(x=1/2\), and
\[
 t_n=\frac{(1/2)^n}{n!\,\Gamma(107/250-(n+1)/3)^2}>0.
\]
Gamma's recurrence gives the exact rational ratio
\[
 \frac{t_{n+3}}{t_n}
 =\frac{(n+4-321/250)^2}{72(n+1)(n+2)(n+3)}
 <\frac{n+3}{72(n+1)(n+2)}.
\]
For \(n\ge30\), the last expression is at most
\(33/71424<1/1000\).  Therefore
\begin{equation}\label{spec:tail-certificate}
 \left|\sum_{n=30}^\infty(-1)^nt_n\right|
 \le\frac{1000}{999}(t_{30}+t_{31}+t_{32})=:T.
\end{equation}
The accompanying script \texttt{leroy\_rational\_certificate.py}
uses exact rational inputs and 256-bit Arb balls for the 33 reciprocal
Gamma evaluations.  It verifies, by interval comparisons,
\[
 T<10^{-30},\qquad
 \Gamma(107/250)^2\left(\sum_{n=0}^{29}(-1)^nt_n-T\right)>-1/2000,
\]
and
\[
 \Gamma(107/250)^2\left(\sum_{n=0}^{29}(-1)^nt_n+T\right)<-1/2500.
\]
For readability, the normalized partial sum is enclosed by the short
decimal interval
\[
 -9.6084728959550860\times10^{-5}
 <\sum_{n=0}^{29}(-1)^nt_n
 <-9.6084728959550858\times10^{-5},
\]
and the absolute normalized remainder bound is below
\(3.838\times10^{-31}\).  The script's rational comparisons, rather
than rounded displayed decimals, are the certificate; its log records
the full Arb enclosures.  Formula
\eqref{spec:tail-certificate} proves \eqref{spec:q-negative};
the second pair is treated identically.  Indeed both pairs satisfy
\(1/3<\beta<1/2\) and \(0<x\le13/25\), so for \(n\ge30\) their
three-step absolute term ratios are at most
\[
 (13/25)^3\frac{33}{9\cdot31\cdot32}
          =\frac{24167}{46500000}<\frac1{1000}.
\]
For the second pair the absolute normalized tail is below
\(7.707\times10^{-31}\); the same script verifies the stated rational
interval, enclosing its value near \(-5.18172\times10^{-13}\).
Finally \eqref{spec:CM-positive} and Theorem~\ref{spec:threshold}
turn each negative-density witness into the corresponding strict lower
bound.  No global uniqueness of a numerical contact is assumed.
\end{proof}

\subsection{A convex-order bridge to the exponential law}

Let \(X_\beta\) have density \(q_\beta\) for \(\beta\ge\beta_c\), and set
\[
 \mu_\beta=\mathbb E X_\beta
       =\frac{\Gamma(\beta)^\gamma}{\Gamma(\beta+\alpha)^\gamma},
 \qquad Y_\beta=X_\beta/\mu_\beta.
\]
We use \(U\le_{\rm cx}V\) to mean that \(U,V\) have the same mean and
\(\mathbb E\varphi(U)\le\mathbb E\varphi(V)\) for every convex function
for which the expectations are defined.

\begin{theorem}[An exponential upper law in convex order]
\label{spec:convex-order}
If \(\beta_c\le\beta_1<\beta_2\), then
\[
 X_{\beta_2}\le_{\rm st}X_{\beta_1},\qquad
 Y_{\beta_1}\le_{\rm cx}Y_{\beta_2}\le_{\rm cx}E,
 \qquad E\sim\operatorname{Exp}(1).
\]
Moreover \(Y_\beta\) converges in distribution to \(E\) as
\(\beta\to\infty\).  The coupling can be chosen so that
\((Y_\beta)_{\beta\ge\beta_c}\) is a positive martingale, bounded in
every fixed \(L^p\), and converges in every finite \(L^p\) to an
exponential terminal variable.
\end{theorem}

\begin{proof}
In \eqref{spec:coupling}, \(0<Z<1\), proving stochastic decrease of
the unnormalized law.  After normalization,
\[
 Y_{\beta_2}\overset d=Y_{\beta_1}V,
 \qquad V=Z/\mathbb E Z,\qquad \mathbb E V=1,
\]
with independence.  Conditional Jensen proves the first convex-order
inequality.

For every positive integer \(n\),
\[
 \mathbb E Y_\beta^n
 =n!\left[
   \frac{\Gamma(\beta+\alpha)^n}
        {\Gamma(\beta)^{n-1}\Gamma(\beta+n\alpha)}
       \right]^\gamma\le n!,
\]
where the inequality is Gamma log-convexity.  The Gamma-ratio
asymptotic makes these moments tend to \(n!\).  The moment problem
for the exponential distribution is determinate; tightness and the
uniform higher-moment bounds prove convergence in distribution to
\(E\).  The second-moment bound supplies uniform integrability,
so the expectations of each stop-loss function \((y-t)_+\) converge.
Passing to the limit in the convex-order inequalities proves
\(Y_\beta\le_{\rm cx}E\).

Finally the logarithms of the factors in \eqref{spec:Z-mellin}
add when consecutive parameter intervals are concatenated.  Independent
increments can be realized using a Poisson measure on the parameter and
jump variables \((b,t)\) with intensity
\[
 \gamma\frac{e^{-bt}}{1-e^{-t}}\,db\,dt,
 \qquad b>\beta_c,\quad t>0,
\]
where a point produces the logarithmic downward jump \(-\alpha t\).
The integral of \(\min(1,t)\) over a finite parameter interval is
finite, so this defines the required multiplicative process.
Dividing it by its mean gives the conditional-mean-one factors above,
and hence a positive martingale \(Y_\beta\).  Uniform bounds in
\(L^n\) for every integer \(n\) give martingale convergence in every
fixed finite \(L^p\); its already identified limiting distribution is
exponential.
\end{proof}

\begin{corollary}[Sharp uniform transform bounds]\label{spec:transform-bounds}
For every \(\beta\ge\beta_c\) and \(x>0\),
\[
 e^{-x}<\Gamma(\beta)^\gamma
             F_\beta(-x/\mu_\beta)<\frac1{1+x}.
\]
For \(0<x<1\),
\[
 e^x<\Gamma(\beta)^\gamma
             F_\beta(x/\mu_\beta)<\frac1{1-x}.
\]
The upper bounds are sharp in the parameter limit \(\beta\to\infty\).
Also \(\mathbb E(Y_\beta-t)_+\le e^{-t}\) for every \(t\ge0\).
\end{corollary}

\begin{proof}
Apply Jensen and Theorem~\ref{spec:convex-order} to \(e^{-xy}\),
\(e^{xy}\), and \((y-t)_+\).  The density is nondegenerate, so Jensen
is strict.  The independent beta factor is nondegenerate for every
positive parameter increment, making each finite-step inequality
strict for the exponential test functions; a further parameter
increment followed by the limit proves strictness of the upper bounds.
For \(0<x<1\), the uniform moment bound by \(n!\) also justifies
convergence of the moment-generating series to \((1-x)^{-1}\).
\end{proof}

\paragraph{What remains open here.}
The results locate a unique true threshold and characterize its finite
moment approximations and local density contacts.  They do not evaluate
\(\beta_c\) in closed form, prove that there is exactly one critical
contact, or prove that every contact is quadratic.  Numerical solutions
of \(q=q_x=0\) are not used to certify the global threshold.  The
strictly negative rational density value in
Theorem~\ref{spec:rational-witness} is a certified lower bound only.

\section{Precise open endpoints suggested by the proofs}
\label{sec:frontiers}

The following are research targets, not additional proved statements.
They are included to show where each bridge reaches and what further
argument would actually be needed.

\paragraph{An individual value in the Erd\H{o}s perturbation class.}
For an arbitrary integer sequence $a_n/2^{2^n}\to1$, the common-field
theorem proves that at most one rational shift can yield a rational
value, and at most two can yield values in a specified quadratic
field. Erd\H{o}s 263 asks for irrationality at the prescribed shift
zero throughout that class \cite{Erdos263}. Moving from ``at most
one exception'' to ``zero is never exceptional'' requires an argument
that locates the exception. The norm argument does not supply it.

\paragraph{Exact evaluation and geometry of the Le Roy boundary.}
Theorem~\ref{spec:threshold} and Corollary~\ref{spec:largest-contact}
reduce the boundary to the largest solution of $q=q_x=0$.
It remains to determine whether the critical density has exactly
one zero and whether each critical zero is always quadratic.
An affirmative nondegeneracy statement would force the
$\varepsilon^{3/2}$ negative-mass law universally. The rigorous
lower bound in Theorem~\ref{spec:rational-witness} gives neither
an equal upper bound nor a closed expression for $\beta_c$.

\paragraph{Effective Bessel selection.}
The explicit $D_r$ family needs no selection procedure. In contrast,
the independent raw $K_0$ and six-coordinate subsets in
Theorem~\ref{arb:rationalheight} are existential. A method that
identifies the exceptional arguments, or settles a specified raw
$K_0(t)$, would go beyond the rank estimate. No irrationality or
transcendence claim for Euler's constant follows from cancelling it.

\paragraph{The binary sign restriction.}
The prime-power construction gives coefficients in a prescribed
congruence class, with exact multiplicity and one minimax solution.
Forcing every off-zero coefficient into $\{-1,1\}$ is an additional
nonlinear restriction. The displayed period-81 and period-64 vectors
do not meet that restriction, so they are not sign-restricted
counterexamples to Erd\H{o}s-type harmonic problems.

\paragraph{Beyond the nearby-prime binomial family.}
The four-value histogram depends on $p<q<r$, increasing gaps, and
$2(r-p)^2<p$. Removing those hypotheses, handling arbitrary squarefree
integers, or proving a global extremal bound for Erd\H{o}s 700 requires
new information. The exact distribution here gives a family on which
such a proposed global result must agree.

\paragraph{From universal to individual integration stability.}
The geometric iff statements concern all regular functions on an
affine curve, or all functions in its algebraic function field.
They give a converse to the polynomial-cover construction and a
dimension for its first obstruction. They do not classify all
individual transcendental elementary functions under repeated
integration, the wider setting of \cite{Chen2022}.

\appendix
\section{A scope check using classical Mahler specialization}
\label{app:mahler}

For completeness, all allowed rational translates of the unperturbed
sequence $2^{2^n}$ are already transcendental by classical Mahler
theory. Here is a direct verification of the needed hypotheses; it
explains the use of the entire perturbation class in
Section~\ref{disc:fieldsection}.

For $h\in\mathbb Q$, set
\[
 g_h(z)=\sum_{n\ge0}\frac{z^{2^n}}{1+h z^{2^n}},\qquad
 g_h(z)-g_h(z^2)=\frac{z}{1+hz}.
\]
This defines an analytic germ at zero and a meromorphic function in
the unit disk. It is never rational. Indeed, if $h\ne0$ and a
rational $g_h$ had finite nonzero pole set $S$, the functional equation
would give
\[
 \{z:z^2\in S\}\subseteq S\cup\{-1/h\},
 \qquad 2|S|\le |S|+1.
\]
If $S$ is empty, $g_h$ is polynomial, contradicting the pole of the
right side. If $S$ has one member, the inclusion forces $S=\{1\}$
and $h=1$. But a pole of order $m$ at one leaves in the difference a
nonzero leading coefficient multiplied by $1-2^{-m}$, whereas
$z/(1+z)$ is regular at one. This too is impossible. When $h=0$,
the same pole-count argument forces a polynomial, and comparison of
degrees in $g(z)-g(z^2)=z$ excludes that possibility.

A nonrational Mahler function is transcendental over
$\overline{\mathbb Q}(z)$, and Nishioka's specialization theorem
transfers the transcendence degree at regular algebraic points
\cite[Theorem~1]{AdamczewskiFaverjon2023}. Apply it to
\[
 \binom{g_h(z)}1=
 \begin{pmatrix}1&z/(1+hz)\\0&1\end{pmatrix}
 \binom{g_h(z^2)}1.
\]
The point $1/4$ is regular exactly when
$h\ne-2^{2^n}$ for every $n\ge1$, the original pole exclusions.
Thus $g_h(1/4)=\sum_{n\ge1}(2^{2^n}+h)^{-1}$ is transcendental.
For a parameter whose initial power-series disk is small, iterate to
a sufficiently small point on this same regular orbit and return by
the functional equation; only rational terms are added. This
classical corollary is not included in the count of 28 statements.

\section{Reproducibility and computational certificates}
\label{app:checks}

The proof statements concern all allowed inputs. The finite checks
listed here audit formulas and examples; they do not replace the
infinite-range arguments. The Le Roy sign bound is different in one
respect: its numerical part is a rigorous interval certificate,
combined with an analytic bound for the infinite remainder.

\begin{center}
\small
\begin{tabular}{@{}p{0.38\linewidth}p{0.55\linewidth}@{}}
\toprule
Check & Result\\
\midrule
Corpus identity & 130 of 130 Git blob hashes matched the snapshot.\\
Quadratic weighted digits & 4,330 exact prime-valuation cases passed.\\
Binomial example & All 556,559 indices in the stated half range matched independent Legendre valuations.\\
Prime 1093 & Multiplicative order and both prime-power residues matched.\\
Period 64 minimum & The algebraic factorization proves height 8 and uniqueness; independent finite enumeration agreed.\\
Le Roy interval witnesses & Both 256-bit Arb certificates passed; each absolute normalized tail is below $10^{-30}$.\\
\bottomrule
\end{tabular}
\end{center}

The following code is embedded in this single TeX source. The first
listing requires Python 3 and \texttt{python-flint}; it was executed
with \texttt{python-flint 0.9.0}. It uses exact rational inputs,
reciprocal Gamma intervals, and comparisons against exact rational
targets. The second listing requires Python 3 and SymPy.
Neither certificate depends on a floating-point root solver.

\subsection{Rigorous Le Roy signs}
\begin{lstlisting}
"""Rigorous non-CM witness for a Le Roy inverse-Mellin density.

Requires python-flint (tested with 0.9.0 / FLINT 3.6.0).
The mathematical proof of the series and its remainder bound is in
spectral_bridges.tex.  This script certifies only the stated rational
point, and does not locate the global positivity threshold.
"""

import math
import flint
from flint import arb, fmpq, ctx


def certify(beta, x, lower_target, upper_target):
    n_terms = 30
    assert fmpq(1, 3) < beta < fmpq(1, 2)
    assert 0 < x <= fmpq(13, 25)

    # All inputs are exact rational numbers, converted to Arb balls.
    terms = [
        arb(x) ** n
        * arb(beta - fmpq(n + 1, 3)).rgamma() ** 2
        / math.factorial(n)
        for n in range(n_terms + 3)
    ]
    partial = sum(
        ((-1) ** n * terms[n] for n in range(n_terms)), arb(0)
    )

    # For n >= 30, beta > 1/3 and x <= 13/25 imply
    # t_(n+3)/t_n <= (13/25)^3 * 33/(9*31*32) < 1/1000.
    assert fmpq(24167, 46500000) < fmpq(1, 1000)
    tail_bound = sum(terms[n_terms : n_terms + 3], arb(0)) * arb(
        fmpq(1000, 999)
    )
    gamma_squared = arb(beta).gamma() ** 2
    lower = gamma_squared * (partial - tail_bound)
    upper = gamma_squared * (partial + tail_bound)

    assert tail_bound < arb(fmpq(1, 10**30))
    assert lower > arb(lower_target)
    assert upper < arb(upper_target)

    print("alpha = 1/3; gamma = 2; beta =", beta, "; x =", x)
    print("Normalized 30-term sum:", partial)
    print("Absolute normalized tail bound:", tail_bound)
    print("q(x) lower enclosure:", lower)
    print("q(x) upper enclosure:", upper)
    print("CERTIFIED:", lower_target, "< q_beta(x) <", upper_target, "< 0")
    print("CONCLUSION: beta_CM(1/3, 2) >", beta)


def main():
    ctx.prec = 256
    print("python-flint version:", flint.__version__)
    print("Arb precision (bits):", ctx.prec)
    certify(fmpq(107, 250), fmpq(1, 2), fmpq(-1, 2000), fmpq(-1, 2500))
    certify(
        fmpq(428292245841, 10**12),
        fmpq(511574841194, 10**12),
        fmpq(-1, 10**12),
        fmpq(-1, 10**13),
    )


if __name__ == "__main__":
    main()
\end{lstlisting}

The observed certified inequalities are
\[
 -1/2000<q_{107/250}(1/2)<-1/2500,
\]
and
\[
 -10^{-12}<
 q_{428292245841/10^{12}}(511574841194/10^{12})
 <-10^{-13}.
\]
The normalized absolute remainder bounds were respectively below
$3.838\cdot10^{-31}$ and $7.707\cdot10^{-31}$.

\subsection{Exact discrete checks}
\begin{lstlisting}
"""Exact supplementary checks for the three discrete/arithmetic bridges.

This is reproducibility material for the accompanying proof dossier.
The all-index theorems have analytic proofs; finite checks do not replace them.
Dependencies: Python 3 and SymPy.
"""

from fractions import Fraction
from itertools import product
from sympy import factorint, n_order


def vp(n, p):
    n = abs(int(n))
    if n == 0:
        raise ValueError("valuation of zero is not needed")
    result = 0
    while n % p == 0:
        n //= p
        result += 1
    return result


def check_quadratic_valuations():
    checked = 0
    for r in range(2, 9):
        modulus = 2**r - 1
        for digits in product((0, 1), repeat=r):
            S = sum(digits[i - 1] * 2**(r - i) for i in range(1, r + 1))
            if S in (0, modulus):
                continue
            y = Fraction(S, modulus)
            denominator = y.denominator
            if int(n_order(2, denominator)) != r:
                continue
            T = sum(i * digits[i - 1] * 2**(r - i) for i in range(1, r + 1))
            U = sum(i*i * digits[i - 1] * 2**(r - i) for i in range(1, r + 1))
            z = Fraction(T, modulus) + Fraction(r*S, modulus**2)
            w = (Fraction(U, modulus) + Fraction(2*r*T, modulus**2)
                 + Fraction(r*r*S*(modulus + 2), modulus**3))
            for prefix_length in range(5):
                value = (prefix_length**2*y + 2*prefix_length*z + w) / 2**prefix_length
                for prime, exponent in factorint(denominator).items():
                    prime, exponent = int(prime), int(exponent)
                    if exponent > vp(r, prime):
                        order = int(n_order(2, prime))
                        lift = vp(2**order - 1, prime)
                        assert vp(value.denominator, prime) == exponent + 2*lift
                        checked += 1
    assert checked == 4330
    return checked


def check_binomial_histogram():
    """Independent of Lucas counting: use factorial valuations at each k."""
    primes = (101, 103, 107)
    n = 101*103*107
    counts = {10403: 0, 10807: 0, 11021: 0, n: 0}
    for k in range(2, n//2 + 1):
        common_divisor = 1
        for prime in primes:
            power, valuation = prime, 0
            while power <= n:
                valuation += n//power - k//power - (n-k)//power
                power *= prime
            if valuation:
                common_divisor *= prime
        assert common_divisor in counts
        counts[common_divisor] += 1
    assert counts == {10403: 1224, 10807: 191, 11021: 116, n: 555028}
    assert sum(counts.values()) == 556559
    return counts


def check_wieferich_example():
    prime = 1093
    assert all(prime % d for d in range(2, 34))
    assert [pow(2, m, prime) for m in (364, 182, 52, 28)] == [1, 1092, 27, 121]
    assert pow(2, 364, prime**2) == 1
    assert pow(2, 364, prime**3) == 1064432260


if __name__ == "__main__":
    print("Quadratic valuation cases:", check_quadratic_valuations())
    print("Binomial histogram:", check_binomial_histogram())
    check_wieferich_example()
    print("All exact checks passed.")
\end{lstlisting}

\subsection{Compact source manifest}

The complete 130-entry source manifest is preserved as comments at
the end of this TeX file. The hashes below identify the principal
inputs; Git blob SHA-1 is used for repository identity, not as a
mathematical validity test.

\begin{center}
\small
\begin{tabular}{@{}ll@{}}
\toprule
Manuscript & Git blob SHA-1\\
\midrule
{[044]} & \texttt{26a390406cc2d22d2cae65b8b67959f8c65b0b9e}\\
{[056]} & \texttt{6a45606058e3d5e5349119fe0c03d981b748d011}\\
{[080]} & \texttt{b86c948fcde18e637bf7144e627005409a7956ea}\\
{[086]} & \texttt{a0c9fb110464b1996c08f1d9d2f11216e98cf862}\\
{[088]} & \texttt{9527133f8f4693075e6ec6f08ec19d46bf29eaa9}\\
{[112]} & \texttt{625d08cfedfe99dff4810de08f9c632e4bb43172}\\
{[113]} & \texttt{8335b4db16c194be3b340f4e06cf4943e5625066}\\
{[114]} & \texttt{33f98aaff53b8512ac115ab2202fedeec83f5756}\\
{[116]} & \texttt{179f647a09a72ed0c9052b2b1d0dc6f82347abe9}\\
\bottomrule
\end{tabular}
\end{center}

\begin{thebibliography}{99}
\raggedright

\bibitem{VMRC}
QEEDDD, \emph{Verifiable Mathematical Reasoning Corpus},
130-manuscript snapshot, commit
\texttt{c30d40119134a59eda6fa5cc6331f27c8e67d034}, accessed
8 November 2026.
\url{https://github.com/QEEDDD/Verifiable-Mathematical-Reasoning-Corpus}.

\bibitem{Beukers2006}
F. Beukers, \emph{A refined version of the Siegel--Shidlovskii theorem},
Annals of Mathematics \textbf{163} (2006), 369--379.
\url{https://doi.org/10.4007/annals.2006.163.369}.

\bibitem{Simon2021}
T. Simon, \emph{Remark on a Mittag-Leffler function of Le Roy type},
arXiv:2012.11018v2 (25 March 2021).
\url{https://arxiv.org/abs/2012.11018}.

\bibitem{DLMF}
NIST, \emph{Digital Library of Mathematical Functions},
Section 5.11, Gamma-function asymptotics.
\url{https://dlmf.nist.gov/5.11}.

\bibitem{Bell2023}
J. P. Bell, K. Gunn, K. D. Nguyen, and J. C. Saunders,
\emph{A general criterion for the P\'olya--Carlson dichotomy and application},
Transactions of the American Mathematical Society \textbf{376} (2023),
4361--4382; arXiv:2206.00862.
\url{https://doi.org/10.1090/tran/8876}.

\bibitem{Chen2022}
S. Chen, \emph{Stability Problems in Symbolic Integration},
ISSAC 2022, 517--524; arXiv:2202.06305.
\url{https://doi.org/10.1145/3476446.3535502}.

\bibitem{Erdos263}
P. Erd\H{o}s and R. L. Graham, \emph{Old and New Problems and Results
in Combinatorial Number Theory}, 1980, p. 63.
Current problem statement and references: T. F. Bloom,
\emph{Erd\H{o}s Problem 263}, consulted 8 November 2026.
\url{https://www.erdosproblems.com/263}.

\bibitem{Coons2011}
M. Coons, \emph{On the rational approximation of the sum of the
reciprocals of the Fermat numbers}, arXiv:1112.5072 (2011).
\url{https://arxiv.org/abs/1112.5072}.

\bibitem{AdamczewskiFaverjon2023}
B. Adamczewski and C. Faverjon,
\emph{A new proof of Nishioka's theorem in Mahler's method},
Comptes Rendus Math\'ematique \textbf{361} (2023), 1011--1028.
\url{https://doi.org/10.5802/crmath.458}.

\bibitem{Zhu2026}
B.-X. Zhu, \emph{Coefficientwise Hankel-total positivity from a
universal cumulant lifting}, arXiv:2611.01172, submitted
31 October 2026.
\url{https://arxiv.org/abs/2611.01172}.

\end{thebibliography}
\end{document}


% Complete audited source manifest (JSON, one entry per line):
% {"commit": "c30d40119134a59eda6fa5cc6331f27c8e67d034", "tex_count": 130, "matching_files": 130}
% {"file": "[001] Exponential-Integral Periods across Arithmetic, Pade, Hankel, Reciprocal Geometry, and Frobenius Framing.tex", "blob_sha1": "1a69e5b0c6a76be7feba8869c8648e006c68a3af", "matches_current_tree": true}
% {"file": "[002] The Polyexponential Trace-Tower Program (Moving  Levels, Faber and Inverse Monodromy, Positive Spectra, and Arithmetic Barriers).tex", "blob_sha1": "92aa1c94ae6bde2047ffb2fd2c9bdffe0398ba1d", "matches_current_tree": true}
% {"file": "[003] Integrated Structures and Boundary Spectra in Sparse Prime Codings.tex", "blob_sha1": "20bb58472842dce122595381f7a68cb5b1d9158c", "matches_current_tree": true}
% {"file": "[004] Prime-Zeta Divisor Tomography Local Zero Processes, Frozen-Edge Rigidity, and Derivative Caustics.tex", "blob_sha1": "7f3c95325da4bb88c11c031a3e2e152c8d174f03", "matches_current_tree": true}
% {"file": "[005] Composite Integer Parts of Powers of Algebraic Units (Pure Cubic Fields, Salem Numbers, and a K3 Application).tex", "blob_sha1": "fc8658e207fa3decd726453a8169ae6998e5670b", "matches_current_tree": true}
% {"file": "[006] Minimal Recurrence Orders and Integral Solution Lattices for Franel Sequences.tex", "blob_sha1": "5e8cb8ae34032f7216a9abd75beb546f45b6c4fe", "matches_current_tree": true}
% {"file": "[007] Dyadic Lucas Prime Cascades and Composite Fibonacci Shifts.tex", "blob_sha1": "11db912f05ecf71d4ba2c36c8c2f7aff0e74423b", "matches_current_tree": true}
% {"file": "[008] Additive Clocks, Non-Reversible Recurrences, and Transcendental Prime-Representing Constants.tex", "blob_sha1": "e84bc6d81245ce037540c2b450758e3b6ff76ad2", "matches_current_tree": true}
% {"file": "[009] Composite Terms in Sparse Cubic Pisot Floors (Alternating Shifts, 3-Adic Limits, and Projective Return).tex", "blob_sha1": "f5a1f066c98e680adbbff0c43045675527833f90", "matches_current_tree": true}
% {"file": "[010] Arithmetic Constraints on a Hypothetically Algebraic Mills Constant (Gap Divisibility, Transcendental Exponents, and Cubic.tex", "blob_sha1": "21eeec56af5b2feec8563ff7f1e1fa5c4ad84382", "matches_current_tree": true}
% {"file": "[011] Rational Perron Bases with Irrational Characteristic Roots (Cubic chamber moments and arithmetic barriers for Catalan approximations).tex", "blob_sha1": "5cc18d9e0af9d20e4efb48a03ecf2fc2b83f9173", "matches_current_tree": true}
% {"file": "[012] Quantitative Word Occurrence in Lambert Constants (Programming Divisor Coefficients in Their Corresponding Radices).tex", "blob_sha1": "82db53ee8ab5b06f52bef298bb958d07b5dfbdd6", "matches_current_tree": true}
% {"file": "[013] PROFINITE RECOGNITION OF MAHLER EQUATIONS AND D-FINITE WITT ESCAPE.tex", "blob_sha1": "6c97b2364c98068a14fa3a666b61708943f81ee9", "matches_current_tree": true}
% {"file": "[014] ENDPOINT LOCALIZATION FOR TYPE III AND TYPE IV SALEM–BOYD INTERLACING FAMILIES.tex", "blob_sha1": "d537fbf861cee55a4de4a80500fc88c1a19aaedd", "matches_current_tree": true}
% {"file": "[015] ONE MISSING EDGE (THE EXACT KN − e CONE, STRICT SUPPORT, AND SHARP DEFECT LAWS).tex", "blob_sha1": "5254dcdeb32304e640d62c437240a93b1dea95cc", "matches_current_tree": true}
% {"file": "[016] UNIVERSAL WEIGHTED-LAPLACIAN CONES AND THE MAXIMAL-CUSP SPECTRAL LINK.tex", "blob_sha1": "cf0a510d0135ab281fd0b015ed5e9c42e912567a", "matches_current_tree": true}
% {"file": "[017] FLOQUET SPARSE OBSTRUCTION (EXACT EXCLUSION THROUGH SUPPORT SIX IN THE PERIOD-THREE HAMMING FIBRE).tex", "blob_sha1": "9541ae4e1dc5e77a548627d280a1ccbf9a20930b", "matches_current_tree": true}
% {"file": "[018] SPECTRAL ARITHMETIC SURGERY PINNED FULL–MINOR ASSIGNMENT, ALL–RANK.tex", "blob_sha1": "3707a7df26a9446ad8034045605ff02e891dfbe2", "matches_current_tree": true}
% {"file": "[019] Euler’s Constant Complete Integrated Research Manuscript (Exponential-Integral Foundations, Positive Moments, Infinite Families, Local Arithmetic, and Transfer Obstructions).tex", "blob_sha1": "e2e470e127d7669bba05756d37e30606275d1d07", "matches_current_tree": true}
% {"file": "[020] The Two-World Search for Euler’s Constant (Shear invariance, collision worlds, forgeries, and the map of possible worlds).tex", "blob_sha1": "c89e2f24340a9d109da43ef5672602b89d71dbc6", "matches_current_tree": true}
% {"file": "[021] Arithmetic Alternatives for Euler’s Constant.tex", "blob_sha1": "3f786ac36a36356de35c5cc8d40fa82ec8323f78", "matches_current_tree": true}
% {"file": "[022] Arithmetic Transcendence of Integrals, Quantiles, and Lemniscates (𝐸-functions, cyclotomic periods, local logarithms, and the obstruction at Euler’s constant ).tex", "blob_sha1": "9c86c9c5fe5aac6236da830689460d0afc294f5f", "matches_current_tree": true}
% {"file": "[023] Explicit Transcendence Criteria for Beta-Type Integrals (Integral Quotients, Cyclotomic Fields, Statistical Quantiles, and Transfers to the Gamma Function).tex", "blob_sha1": "75c2e3a820b0e049e1819b546a003b8e01e8807f", "matches_current_tree": true}
% {"file": "[024] Transcendence of Algebraic Arc Lengths and Inverse Periods (Division Dichotomies, Beta Quantiles, and Arithmetic Transfers).tex", "blob_sha1": "873704debce016d08bed812eccff996362bfb564", "matches_current_tree": true}
% {"file": "[025] Gamma-Type Functions (Arithmetic and Differential Independence, Probability Laws, and Interpolation).tex", "blob_sha1": "d325b86ddcb8ff7330c16afb9b26be5c824ef153", "matches_current_tree": true}
% {"file": "[026] Exact Orders, Shift Optimality, and Zero Asymptotics for q-Exponential Hankel Determinants.tex", "blob_sha1": "443e62b866821fc9c289d030454d71145a418859", "matches_current_tree": true}
% {"file": "[027] Arithmetic of Ball Overlaps and Their Iterated Volume Maps.tex", "blob_sha1": "db0ae51bc2e428ac64e8c3d38b48fb7318698dd9", "matches_current_tree": true}
% {"file": "[028] Observation and Recovery in Arithmetic and Probability (Prime Flows, Rare Events, and Infinite Mixtures).tex", "blob_sha1": "8ffba42d2165a4d8b314dcfb9cf5281cd0e79807", "matches_current_tree": true}
% {"file": "[029] Observation and Recovery (Geometry, Mixtures, and Spectral Inversion).tex", "blob_sha1": "336b9be8f2c9d47f8b26a5b50edb201acf8bbfdb", "matches_current_tree": true}
% {"file": "[030] Arithmetic and Observation Classifications, Recovery, and Transcendence Barriers.tex", "blob_sha1": "941b88336ccbad3c2f1553ac2d09202576ada8d8", "matches_current_tree": true}
% {"file": "[031] Prime-Representing Constants and Related Structures (Encodings, Arithmetic, Geometry, and Modular Obstructions).tex", "blob_sha1": "da444ea874b06d10ce8ed1110cca4da8725ed619", "matches_current_tree": true}
% {"file": "[032] Minimal Differential and Recurrence Orders (Bessel Periods, Bounded-Height Tableaux, and Banana Integrals).tex", "blob_sha1": "b926f675eb65117b85d32a515539e34ea7ed13e7", "matches_current_tree": true}
% {"file": "[033] Euler’s Constant Arithmetic, Monodromy, and Gamma-Related Structures.tex", "blob_sha1": "58f9b37707147be1c90bae67769c9063e6e68b8b", "matches_current_tree": true}
% {"file": "[034] Global Proof Transfers.tex", "blob_sha1": "7b6abe250b10b15d762e61fd0c9c60af617e438c", "matches_current_tree": true}
% {"file": "[035] Proof Transfers and Global Classifications (Complete Integrated Edition Spectral Design, Statistical Recovery, Recurrences, Period Geometry, Arithmetic, and Exact Cone Gluing).tex", "blob_sha1": "1a017687873c0c4cb4a907b8da5ce23303932aa3", "matches_current_tree": true}
% {"file": "[036] Exact Recurrences, Algebraic Independence, and Certified Statistical Recovery (Vertex balance, Bessel modules, and entropy transport).tex", "blob_sha1": "c0f4a42843922ba808a0d42409e8acecaca06f79", "matches_current_tree": true}
% {"file": "[037] ISOTROPIC FUCHS SQUEEZES AND EXACT DIFFERENTIAL COMPLEXITY (HYPERCUBIC GREEN FUNCTIONS, REGEV SUMS, AND RANDOM WALK POLES).tex", "blob_sha1": "ae24434aa2da76f2ee53e99de05323765a772d3f", "matches_current_tree": true}
% {"file": "[038] Arithmetic Rank and Differential Complexity of Hamiltonian Moment Integrals.tex", "blob_sha1": "a17681a076b6529d212299b1b5856398ce715366", "matches_current_tree": true}
% {"file": "[039] Temporal Sampling and Local Nonidentifiability in Coalescent Demography.tex", "blob_sha1": "a1b55d1c380a9f61a469806cfc09104cc50cb373", "matches_current_tree": true}
% {"file": "[040] OEIS Conjecture Hunt.tex", "blob_sha1": "2587c93cb1aa02237cff6229718fff63b96b4f46", "matches_current_tree": true}
% {"file": "[041] Ewald points of monotone polytopes the minimum in dimension eight and branching simplex bundles.tex", "blob_sha1": "c1cf25dbd56c8016ba7707fc7da1afc9dc7ed9a9", "matches_current_tree": true}
% {"file": "[042] Cubic Pisot Floors (Composite Windows, Sparse Exceptions, and Finite-Sieve Barriers).tex", "blob_sha1": "586b9087bd5bdc229bda403c6d068d0bd7ce022c", "matches_current_tree": true}
% {"file": "[043] Local Certificates, Stability, and Concavity.tex", "blob_sha1": "714975ec9c589882717b657b9e619199d02ff408", "matches_current_tree": true}
% {"file": "[044] Euler–Gompertz and Bessel Values (Joint Independence, Integral Intersections, and Critical Limits).tex", "blob_sha1": "26a390406cc2d22d2cae65b8b67959f8c65b0b9e", "matches_current_tree": true}
% {"file": "[045] Arithmetic of Gieseking Moment Lattices (Normality, Integral Saturation, and Approximation Budgets).tex", "blob_sha1": "40271e33448645ee914fb5c2be9929c117fedf30", "matches_current_tree": true}
% {"file": "[046] Cross-field Consequences of Bessel-Value Independence (Circular Mixtures, Random-Time Quantum Walks, and Robin Spectra).tex", "blob_sha1": "1981018c525499db9309c6a1fb480ce88d28b4dc", "matches_current_tree": true}
% {"file": "[047] A monic resultant for Bessel derivative coefficients.tex", "blob_sha1": "6d049e577dc677ef93e05d8aba32993e6d00501e", "matches_current_tree": true}
% {"file": "[048] Recurrences, Absolute Differences, and Collatz Orbits (Algebra, Prime Patterns, and Zero Distributions).tex", "blob_sha1": "5d262c7cdcf61493cea7be18a255b6be9200e810", "matches_current_tree": true}
% {"file": "[049] Winding Embeddings and Prime-Power Pisot Floors (Integral Intersections, Composite Windows, and Local Obstructions).tex", "blob_sha1": "4bf8405ccbb8c097435655b4936afdf2c087f503", "matches_current_tree": true}
% {"file": "[050] Exact Prime Covers and Arithmetic Obstructions (Salem powers, repeated blocks, Euler integrals and Lambert functions).tex", "blob_sha1": "6b266dc57fc9c3f1c5f8ef207615b399e46e3784", "matches_current_tree": true}
% {"file": "[051] A Proof of Colton’s Conjecture on Refactorable Numbers Smooth–rough τ (decomposition, certified checkpoint chains, and a template tail).tex", "blob_sha1": "9653ed5135f2aaa296fd806e8d578e84226f75ce", "matches_current_tree": true}
% {"file": "[052] Rooted Induced Trees and the Radius_A Proof of Conjecture 84 of Written on the Wall II.tex", "blob_sha1": "49d06a8cc337aa826c0f4980e8daef4213b5d2d3", "matches_current_tree": true}
% {"file": "[053] RESTRICTION CLASSES AND BAKER’S QUESTION ON THE LINEAR INDEPENDENCE OF L(1, χ).tex", "blob_sha1": "ee2a75cf712b536df5ac5d6aa45e2dd5d6c11d34", "matches_current_tree": true}
% {"file": "[054] Supportive Edges, Radius-Stable Peeling and Rooted Yang–Lu Replacement_the Randić–Radius Conjecture for Graphs of Cyclomatic Number at most 8.tex", "blob_sha1": "f9f3a09cc60f490cd1691270b8596c43b8bb68d5", "matches_current_tree": true}
% {"file": "[055] ERDOS CONJECTURE.tex", "blob_sha1": "5b381385a365a19b510a25725f0d3b078134af56", "matches_current_tree": true}
% {"file": "[056] SQUARE MODULI, SPARSE MASKS, AND SIGN OBSTRUCTIONS FOR VANISHING ERDOS SUMS .tex", "blob_sha1": "6a45606058e3d5e5349119fe0c03d981b748d011", "matches_current_tree": true}
% {"file": "[057] erdos_nonsquare_moduli.tex", "blob_sha1": "64b833115332720bb714af6644051812254a3c68", "matches_current_tree": true}
% {"file": "[058] Exact_Arithmetic_at_Small_Salem_Numbers.tex", "blob_sha1": "efe9deefa2f8e32bbdd9b34c8f14f6b4aa3c57d4", "matches_current_tree": true}
% {"file": "[059] Integrated_Mathematical_Research_60.tex", "blob_sha1": "6d00efba2a3a1e1f855da9848647b36b5f4754c7", "matches_current_tree": true}
% {"file": "[060] Corvaja_Zannier.tex", "blob_sha1": "ebde16141269a4937f0937a4fa138ae46b83674e", "matches_current_tree": true}
% {"file": "[061] floquet_level.tex", "blob_sha1": "6a83cdc58b38121a03a1acea11fbef28bfa78425", "matches_current_tree": true}
% {"file": "[062] Gamma_Derivatives_Unified_Windows_and_Counts.tex", "blob_sha1": "6b8fcae7c6f67c0b3b401e48fcd26df1e9fa9bc3", "matches_current_tree": true}
% {"file": "[063] Stanley_Depth.tex", "blob_sha1": "e0fb4ab6432470036643ac72a514436558d699c8", "matches_current_tree": true}
% {"file": "[064] ternary_trace_dynamics_shifted_Mills.tex", "blob_sha1": "6a7d6290c71cb93d750762df8013ea17387a6430", "matches_current_tree": true}
% {"file": "[065] WOWII_Unified.tex", "blob_sha1": "cbf159fe07b49b73e3b96b36f6c71a7f7b142d2e", "matches_current_tree": true}
% {"file": "[066] WOWII_Rounds_01_22_Consolidated.tex", "blob_sha1": "e933a61e830f0fa1e57d75e56474571f03cb6183", "matches_current_tree": true}
% {"file": "[067] open_problems_transferred.tex", "blob_sha1": "7c7036ba787b6a03f98a63b6d2a1f9a779ab500f", "matches_current_tree": true}
% {"file": "[068] orbit_constraints.tex", "blob_sha1": "64cbeaf78b260d4dfc14f1570ce41adff57760a5", "matches_current_tree": true}
% {"file": "[069] Antisymmetric_Spectral_Invariants.tex", "blob_sha1": "d46394b864713e554e022a7fd12fb43f4063b517", "matches_current_tree": true}
% {"file": "[070] Frobenius_Rigidity_Stabilization.tex", "blob_sha1": "4be884bc930a37beb11e753275e25bbcfd06e81b", "matches_current_tree": true}
% {"file": "[071] Obstructions_Polytopes_Peak_Transfer.tex", "blob_sha1": "9eb017a4aa33016ca9b4579c3170eb77dc43b4de", "matches_current_tree": true}
% {"file": "[072] unified_collatz_recurrences.tex", "blob_sha1": "5d262c7cdcf61493cea7be18a255b6be9200e810", "matches_current_tree": true}
% {"file": "[073] Free_Spider_Counterexample.tex", "blob_sha1": "02566f6811502b05ede2bb1c426cae747bcbec58", "matches_current_tree": true}
% {"file": "[074] Collected_Research.tex", "blob_sha1": "3a8ff2947b4259520b3e4fa8bc68f321f48b844b", "matches_current_tree": true}
% {"file": "[075] Shifted_Mills_and_Cubic_Pisot_Traces.tex", "blob_sha1": "854a1c99c2c90b303df20a20b66f70ca773ce42d", "matches_current_tree": true}
% {"file": "[076] T91_T95_Retry.tex", "blob_sha1": "2d6b5a6cb58a73c6194650aac44eaf93b2e4fccd", "matches_current_tree": true}
% {"file": "[077] T96_T100 .tex", "blob_sha1": "c90c1609d9d19cffa9912aa215a7774c4aa4b052", "matches_current_tree": true}
% {"file": "[078] mixed_mahler_congruences.tex", "blob_sha1": "c61fbd4fafa3281840f9d71fab59f8def303e2d8", "matches_current_tree": true}
% {"file": "[079] T86_T90_Stern_Roots_and_Hankel_Recurrences.tex", "blob_sha1": "f7b19d2b2bd2f3b332e458d863bd8af15e17e6b4", "matches_current_tree": true}
% {"file": "[080]_exponent_window_erdos265.tex", "blob_sha1": "b86c948fcde18e637bf7144e627005409a7956ea", "matches_current_tree": true}
% {"file": "[081]_carry_forcing_trinomial_bases.tex", "blob_sha1": "2e0d9b77f5b39766536ee1f7390267ccc3f01cde", "matches_current_tree": true}
% {"file": "[082] A252867_permutation.tex", "blob_sha1": "b189991228ebfdbe9d25524bd92060ded4ccb6a7", "matches_current_tree": true}
% {"file": "[083] beta_curvature_paper.tex", "blob_sha1": "98cbc4a50170efadd3ef14d2b2b7ee6e9f3c2cbe", "matches_current_tree": true}
% {"file": "[084] caporossi_hansen_randic_radius.tex", "blob_sha1": "f9f3a09cc60f490cd1691270b8596c43b8bb68d5", "matches_current_tree": true}
% {"file": "[085] decoder_trees_5093_5107.tex", "blob_sha1": "c90188885d287c389668764273c46b2e610a8fda", "matches_current_tree": true}
% {"file": "[086] erdos700_three_nearby_primes.tex", "blob_sha1": "a0c9fb110464b1996c08f1d9d2f11216e98cf862", "matches_current_tree": true}
% {"file": "[087] fajtlowicz_radius_core_reduction.tex", "blob_sha1": "f969081b3dceffa08355f52059d1aebecd1724ca", "matches_current_tree": true}
% {"file": "[088] Integer_Tail_Analytic_Attack.tex", "blob_sha1": "9527133f8f4693075e6ec6f08ec19d46bf29eaa9", "matches_current_tree": true}
% {"file": "[089] Congruence_Compression_2026.tex", "blob_sha1": "39fbf7631948cfd85d551e600091a07aa68f961c", "matches_current_tree": true}
% {"file": "[090] erdos_graham_tent.tex", "blob_sha1": "55a2edfa803726c7e83b37e1b70e3fd0a96e8fa0", "matches_current_tree": true}
% {"file": "[091] symmetric_insertion_Z.tex", "blob_sha1": "8608abd73d49165a24ac06ea38016f7364ef3e98", "matches_current_tree": true}
% {"file": "[092] kesten_fixed_start.tex", "blob_sha1": "310732f57a3642655e3d48930215840175f725db", "matches_current_tree": true}
% {"file": "[093] pm1_vanishing_paper.tex", "blob_sha1": "17e7a423ae45f9160abcd552bbf28aea5ff58638", "matches_current_tree": true}
% {"file": "[094] Randic_radius_bag_reduction.tex", "blob_sha1": "2f215796b32c8201d16b10903d2d8bd9f9e02280", "matches_current_tree": true}
% {"file": "[095] Simultaneous_Fiber_Cancellation (1).tex", "blob_sha1": "23d078cc90527b883469594fbc70d770283ca233", "matches_current_tree": true}
% {"file": "[096] arctangent_inverse_complete_monotonicity.tex", "blob_sha1": "eba1c3856741b7d87e70e354e90531e7790c4b3c", "matches_current_tree": true}
% {"file": "[097] dirichlet_completeness_criterion.tex", "blob_sha1": "dde13dbb94a674af3900c4ff87af3d7206e86fcb", "matches_current_tree": true}
% {"file": "[098] gamma_positive_certificates_primitive_arithmetic.tex", "blob_sha1": "6464d1de132837a4e834d05f90bfc6c8ae7d89d0", "matches_current_tree": true}
% {"file": "[099] inverse_arctan_stieltjes.tex", "blob_sha1": "3857753333dce27696e925785d12d1fa1581fb6b", "matches_current_tree": true}
% {"file": "[100] Sondow_Grid_Research_Checkpoint.tex", "blob_sha1": "18e800d62c676667af4780acd8bde1ace6ed0979", "matches_current_tree": true}
% {"file": "[101] robin_convex_domain_nonconcavity.tex", "blob_sha1": "45016f5634fbbabdd138f8cddb7453c156bb7094", "matches_current_tree": true}
% {"file": "[102] robin_dilation_nonconvexity.tex", "blob_sha1": "9df78f87a29b65b1a4f84857e517f934d1d8d1f1", "matches_current_tree": true}
% {"file": "[103] A_Nodal_Criterion_for_Real_Zeros.tex", "blob_sha1": "06cb7cc92b7051580b469e41546f80e64c3a2af3", "matches_current_tree": true}
% {"file": "[104] Fixed_Shifted_Mills_Minima.tex", "blob_sha1": "58312a29ce634095b21348946ca1caaf2f1e832b", "matches_current_tree": true}
% {"file": "[105] gamma_consolidated_positive_kernels_and_arithmetic_2026_11_02.tex", "blob_sha1": "42642d5aa483cd69dc2e6c7ab33f5d4311b49185", "matches_current_tree": true}
% {"file": "[106] Degree_Seven_and_Quantum_Path_Families_2026_11_02.tex", "blob_sha1": "a79f34a3165b883cefbdcd76e695ef657f74a375", "matches_current_tree": true}
% {"file": "[107] Dyadic_Saturation_and_Companion_Results_2026_11_02.tex", "blob_sha1": "71a944ace908ed0fbeab9e24217e3e7d465eb7fc", "matches_current_tree": true}
% {"file": "[108] Exact_Rook_Recurrence_Certificates_2026_11_02.tex", "blob_sha1": "591be72f1f3a85c0a82903171af2bcca002add75", "matches_current_tree": true}
% {"file": "[109] Explicit_Order_Degree_Curves_for_Dfinite_Telescoping_2026_11_02.tex", "blob_sha1": "6296ef1874b1cef06ec3c7e4f18d8b92f8b8c8d0", "matches_current_tree": true}
% {"file": "[110] Gamma_Mills_Integrated_2026-11-05.tex", "blob_sha1": "43c5b5b8651e32c6d499f5feaea990be8dccb074", "matches_current_tree": true}
% {"file": "[111] Grouped_Diagonals_and_Arithmetic_Telescoper_Bounds_2026_11_02.tex", "blob_sha1": "4278eb13cb18d0f05946ad1dd1235df1912d6a09", "matches_current_tree": true}
% {"file": "[112] Period_Obstructions_for_Weighted_Binary_Expansions_2026_11_02.tex", "blob_sha1": "625d08cfedfe99dff4810de08f9c632e4bb43172", "matches_current_tree": true}
% {"file": "[113] Polynomial_Inverses_Counterexamples_to_Chen_Stability_2026_11_02.tex", "blob_sha1": "8335b4db16c194be3b340f4e06cf4943e5625066", "matches_current_tree": true}
% {"file": "[114] Research_Results_2026-11-02.tex", "blob_sha1": "33f98aaff53b8512ac115ab2202fedeec83f5756", "matches_current_tree": true}
% {"file": "[115] Ten_External_Problems_2026-11-05.tex", "blob_sha1": "cdd0e8b6be566c72538c9b320352415cac6026f8", "matches_current_tree": true}
% {"file": "[116] Le_Roy_Positivity_Beyond_Infinite_Divisibility.tex", "blob_sha1": "179f647a09a72ed0c9052b2b1d0dc6f82347abe9", "matches_current_tree": true}
% {"file": "[117] Mathematical_Research_Manuscript.tex", "blob_sha1": "8cca4b8fc894859f64633b226e1c6389a237243a", "matches_current_tree": true}
% {"file": "[118] Mathematical_Research_Manuscript (1).tex", "blob_sha1": "8cca4b8fc894859f64633b226e1c6389a237243a", "matches_current_tree": true}
% {"file": "[119] Mathematical_Research_Manuscript (2).tex", "blob_sha1": "5b4e61a6f5a996b8bc15e1d7a2a03966fa25f23a", "matches_current_tree": true}
% {"file": "[120] Mathematical_Research_Manuscript (3).tex", "blob_sha1": "5b4e61a6f5a996b8bc15e1d7a2a03966fa25f23a", "matches_current_tree": true}
% {"file": "[121] Mathematical_Research_Manuscript (4).tex", "blob_sha1": "5b4e61a6f5a996b8bc15e1d7a2a03966fa25f23a", "matches_current_tree": true}
% {"file": "[122] Mean_Median_Mode_and_Gamma_Mixtures_of_Stable_Laws.tex", "blob_sha1": "c3b703923904290c5714da708226233d6b8b820e", "matches_current_tree": true}
% {"file": "[123] Neighborhoods_Kernels_and_Fractional_Zeros.tex", "blob_sha1": "bc5f0aca1c317173ab43eb309aece88a7237d531", "matches_current_tree": true}
% {"file": "[124] Prime_Neighbours_of_Shifted_Mills_Outputs.tex", "blob_sha1": "944fea225bce7e1b27cead74fc35efd46fbabee9", "matches_current_tree": true}
% {"file": "[125] Prime_Outputs_and_Local_Counts.tex", "blob_sha1": "f337c8000d954b7004bc8683dabfcd5e3aef61f7", "matches_current_tree": true}
% {"file": "[126] Real_Collisions_in_the_Gaussian_Heat_Flow.tex", "blob_sha1": "4a54065fb4388e841882ea73f847348c4f401cab", "matches_current_tree": true}
% {"file": "[127] Restricted_Partitions_Domination_and_Kanter_Factors.tex", "blob_sha1": "9f13b57fa57230b0047babeebeee4d05e099771e", "matches_current_tree": true}
% {"file": "[128] Shifted_Mills_Minima.tex", "blob_sha1": "4eba1d00804f2c53cc5bece049d5f03be338fe7b", "matches_current_tree": true}
% {"file": "[129] Certified_Prime_Lifting.tex", "blob_sha1": "28dd3632367c37131e05e1b5c3efe97b98bc229b", "matches_current_tree": true}
% {"file": "[130] Certified_Prime_Lifting_and_Sparse_Modular_Squaring.tex", "blob_sha1": "becfdd1daf545191efd7730b3097ad32b6950fd0", "matches_current_tree": true}
